\documentclass[10pt,a4paper]{article}

% ----------------- Paquetes -------------------%
\usepackage[T1]{fontenc}
\usepackage[utf8]{inputenc}
\usepackage[a4paper,margin=2.5cm]{geometry}
\usepackage{mathptmx}             
\usepackage{changepage} 
\usepackage{setspace}
\usepackage[spanish]{babel}
\usepackage[labelfont=bf,labelsep=space]{caption}
\usepackage{amssymb}
\usepackage{amsmath}
\usepackage{ebgaramond}
\usepackage{placeins}
\usepackage{etoolbox}
\usepackage{setspace}
\usepackage{parskip} 
\usepackage{tcolorbox}
\usepackage{needspace}
\usepackage{xparse}
\usepackage{ocgx2}
\usepackage{hyperref}
\usepackage{hypcap}
\usepackage{soul}\sethlcolor{yellow}% resaltado amarillo: \hl{texto} (Panel TCEE, ⌃H)



%%%%%%%%%%%%%%%%%%%%%%%%%%%%%%%%%%%%%%%%%%%%%%%%%%%%%%%%%%%%%%%%
% --- Fija primero tus valores globales "buenos" ---
\setstretch{1.2}                 % Interlineado global
\setlength{\parskip}{0.7\baselineskip}
\setlength{\parindent}{0pt}
\newlength{\GlobalParskip}
\setlength{\GlobalParskip}{\parskip}

\newcounter{eqnum}[subsection]
\renewcommand{\theeqnum}{\arabic{eqnum}}
\newcounter{alphaitem}
\newcounter{numitem}

%% ------------------- Recuadros----------------------%%
\newcommand{\puntosclave}[1]{%
  \begin{tcolorbox}[
    colframe=black,
    title={Puntos clave},
    before upper={\setlength{\parindent}{0.5cm}\setlength{\parskip}{0.5\baselineskip}}
  ]
  #1
  \end{tcolorbox}
}

\newcommand{\modificaciones}[1]{
  \begin{tcolorbox}[
    colback=white,
    colframe=red,
    title={Modificaciones a realizar},
    before upper={\setlength{\parindent}{0.5cm}\setlength{\parskip}{0.5\baselineskip}}
  ]
  #1
  \end{tcolorbox}
}

%% ------------------- Preguntas TEST----------------------%%
% Opciones: radio (single-choice)
\NewDocumentCommand{\OptRadio}{m m m}{%
  \noindent\CheckBox[radio,name=#1,value=#2,width=1.2ex,height=1.2ex]{}%
  \;\textbf{#2.}~#3\par
}

% Opciones: checkbox (multi-choice)
\NewDocumentCommand{\OptCheck}{m m}{%
  \noindent\CheckBox[name=#1,width=1.2ex,height=1.2ex,bordercolor={0 0 0}]{}%
  \;#2\par
}

% Pregunta single-choice (radio)
% Args: 1=id, 2=enunciado, 3=opciones, 4=solución+justificación, 5=imagen (o {}), 6=origen
\NewDocumentCommand{\PreguntaTestSingle}{m m m m m m}{%
  \par\medskip
  \noindent\textbf{#1.}~#2\par\smallskip

    % Imagen (si no es {})
    \IfBlankTF{#5}{}{%
      \begin{center}
        \includegraphics[width=0.75\linewidth]{#5}
      \end{center}
    }

    \begin{Form}
      #3
    \end{Form}

    \par\smallskip
    \switchocg{sol-#1}{\fbox{Mostrar / ocultar solución}}%
    \hfill{\footnotesize #6}\par\smallskip

    \begin{ocg}{Solución}{sol-#1}{0}
      \noindent #4
    \end{ocg}
  \par\medskip
}

% Pregunta multi-choice (checkboxes)
\NewDocumentCommand{\PreguntaTestMulti}{m m m m m m}{%
  \par\medskip
  \noindent\textbf{#1.}~#2\par\smallskip

    % Imagen (si no es {})
    \IfBlankTF{#5}{}{%
      \begin{center}
        \includegraphics[width=0.75\linewidth]{#5}
      \end{center}
    }
    \begin{Form}
      #3
    \end{Form}

    \par\smallskip
    \switchocg{sol-#1}{\fbox{Mostrar / ocultar solución}}%
    \hfill{\footnotesize #6}\par\smallskip

    \begin{ocg}{Solución}{sol-#1}{0}
      \noindent #4
    \end{ocg}
  \par\medskip
}

%% ------------------- Listas----------------------%%
    % --- Lista alfabética ---
    \newenvironment{la}
      {\begin{list}{\alph{alphaitem})}{%
          \usecounter{alphaitem}%
          \setlength{\leftmargin}{1cm}%
          \setlength{\parskip}{3pt}% 
          \setlength{\itemsep}{0pt}% ← espacio entre ítems
          \setlength{\parsep}{0pt}%  ← párrafos dentro de un ítem
      }}
      {\end{list}}
      
    % --- Lista numérica ---
    \newenvironment{lnum}
      {\begin{list}{\arabic{numitem}.}{%
          \usecounter{numitem}%
          \setlength{\leftmargin}{1cm}%
          \setlength{\parskip}{3pt}% 
          \setlength{\itemsep}{0pt}% ← espacio entre ítems
          \setlength{\parsep}{0pt}%  ← párrafos dentro de un ítem
      }}
      {\end{list}}

%% ------------------- Preguntas Test ----------------------%%
      \usepackage{xparse} % si ya lo tienes, no lo repitas

    \NewDocumentEnvironment{test}{m m o}{%
      \begin{tcolorbox}[
        colback=white,
        colframe=black!15,
        boxrule=0.6pt,
        arc=2mm,
        left=2mm,right=2mm,top=2mm,bottom=2mm
      ]
      \centering
      \includegraphics[width=\linewidth]{#1}
      \end{tcolorbox}
    
      \vspace{2mm}
    
      % Caja SOLO para texto (SÍ partible y mantiene márgenes al partir)
      \begin{tcolorbox}[
        enhanced,
        breakable,
        colback=black!3,
        colframe=black!25,
        boxrule=0.6pt,
        arc=2mm,
        left=2mm,right=2mm,top=1.5mm,bottom=1.5mm
      ]
      \textbf{Respuesta:} \textbf{#2}%
      \IfValueT{#3}{\par\smallskip \textbf{Justificación:} #3}
      \end{tcolorbox}
    }{}
      
% ------------------- Imágenes----------------------%
    \graphicspath{{Img/}}% Rutas por defecto 
    % --- Imagen adaptativa ---
    % Registros de dimensión definidos una sola vez
    \newdimen\imgcap
    \newdimen\imgmin
    \newdimen\imgmax
    \newdimen\imgremain
    \newdimen\imgh
    \newdimen\imgneed
    
    \newcommand{\imagenfit}[3]{%
      \addvspace{10pt}% separación antes
      \par\begingroup
        % Parámetros de composición (asignaciones, no \newdimen)
        \imgcap=1\baselineskip            % alto estimado de título+fuente
        \imgmin=0.15\textheight
        \imgmax=0.15\textheight
        \imgremain=\pagegoal
        \ifdim\pagegoal=\maxdimen \imgremain=0pt\fi
        \advance\imgremain by -\pagetotal
        \advance\imgremain by -\imgcap
        % Decisión de altura destino
        \ifdim\imgremain<\imgmin
          \newpage
          \imgh=0.20\textheight
        \else
          \imgh=\imgremain
          \ifdim\imgh>\imgmax \imgh=\imgmax \fi
        \fi
        % Reservar espacio para evitar cortes del bloque completo
        \imgneed=\imgh \advance\imgneed by \imgcap
        \Needspace{\imgneed}
        % Cuerpo
        \noindent\begin{minipage}{\textwidth}
          \centering
          \captionof{figure}{\textemdash\ #2}\par\vspace{0.3em}% título arriba
          \includegraphics[width=\linewidth,height=\imgh,keepaspectratio]{\detokenize{#1}}\par
          \vspace{0.3em}\small\textit{Fuente: }#3% fuente abajo
        \end{minipage}%
      \endgroup
      \par\addvspace{10pt}%
    }

%------------ Bloque de RECUADRO ESPECIAL -------------%
    \newenvironment{specialblock}[1]{%
      \begin{quote} % crea sangría a izquierda y derecha
      \small % texto más pequeño
      \noindent\textbf{#1}\\[-0.3em] % título en negrita
      \rule{\linewidth}{0.4pt}\\[0.6em]}{
      \end{quote}}
      
%-------------------------- Portada -----------------------%
\newcommand{\oldtitle}[1]{\def\OldTitle{#1}}
\newcommand{\changerationale}[1]{\def\ChangeWhy{#1}}

\newcommand{\changebox}{%
\begin{tcolorbox}[title={Historial de cambios del tema}]
\textbf{Título anterior:} \OldTitle\par
\textbf{Motivo del cambio:} \ChangeWhy
\end{tcolorbox}
}

%-------------------------- Author Font -----------------------%
    \newcommand{\authorfont}[1]{{\fontfamily{EBGaramond-TLF}\selectfont #1}}

%-------------------------- Anexos -----------------------%
    \makeatletter
    \newcounter{anexo}
    \renewcommand{\theanexo}{\Roman{anexo}}
    \newcommand{\anexo}[1]{%
      \clearpage
      \phantomsection
      \refstepcounter{anexo}%
      \noindent\textbf{Anexo \theanexo.- }#1\par\medskip
      \addcontentsline{stc}{section}{Anexo \theanexo.- #1}%
    }
    \makeatother

    %-------------------------- Lista de figuras (personal) -----------------------%
\makeatletter
\newcommand{\MyListOfFigures}{%
  \clearpage
  \phantomsection
  {\centering\bfseries\Large Índice de figuras\par}%
  \medskip
  % Entrada en nuestro índice de contenido (stc)
  \addcontentsline{stc}{section}{Índice de figuras}%
  % Recuperación del listado (sin cabecera propia)
  \begingroup
    \parindent=0pt\parskip=2pt
    \@starttoc{lof}%
  \endgroup
}
\makeatother

%-------------------------- Bibliografía -----------------------%
\makeatletter
% Contador para las referencias
\newcounter{myref}
% Cabecera de la bibliografía, entra en el índice 'stc'
\newcommand{\MyBibliography}{%
  \clearpage
  \phantomsection
  {\centering\bfseries\Large Bibliografía y referencias\par}%
  \medskip
  \addcontentsline{stc}{section}{Bibliografía y referencias}%
  \setcounter{myref}{0}%
}
\newcommand{\MyBibItem}[4]{%
  \refstepcounter{myref}%
  \noindent[\themyref]\ #1.\ \emph{#2}. #3. #4\par
}
\makeatother

%--------- Establecer cuatro niveles de índice --------%
    \setcounter{tocdepth}{4}   
    \setcounter{secnumdepth}{4}% numera hasta \paragraph

%%%%%%%%%%%%%%%%%%%%%%%%%%%%%%%%

\makeatletter

    %--------- Interlineado Global --------%
    \edef\GlobalBaseLineStretch{\baselinestretch}
    
    %--------- Espaciado de seccinones --------%
    \renewcommand\section{\@startsection{section}
      {1}{\z@}
      {2.5ex \@plus 1ex \@minus .2ex} % espacio antes
      {1.5ex \@plus .2ex}             % espacio después
      {\normalfont\Large\bfseries}    % formato del título
    }
    \renewcommand\subsection{\@startsection{subsection}{2}{\z@}
      {3ex \@plus 0.8ex \@minus .2ex}
      {1.5ex \@plus .2ex}
      {\normalfont\large\bfseries}
    }
    \renewcommand\subsubsection{\@startsection{subsubsection}{3}{\z@}
      {3ex \@plus 0.5ex \@minus .2ex}
      {1ex \@plus .2ex}
      {\normalfont\normalsize\bfseries}
    }
    \renewcommand\paragraph{\@startsection{paragraph}{4}{\z@}%
    {-2ex \@plus -0.5ex \@minus -.2ex}% espacio antes
    {0.8ex \@plus .2ex}% espacio después
    {\normalfont\normalsize\itshape}}% ← cursiva en lugar de negrita

    %--------- Ecuaciones --------%   
    % Variables globales para almacenar títulos y notas
    \gdef\leq@current@section{}%
    \gdef\leq@current@subsection{}%
    \gdef\leq@last@printed@section{}%
    \gdef\leq@last@printed@subsection{}%
    
    % Comando para añadir nota (invisible en el cuerpo)
    \newcommand{\eqnote}[1]{%
      \addcontentsline{leq}{leqnote}{#1}%
    }
    
    % Comando para ecuaciones
    \newcommand{\eqblock}[2]{%
      % Verificar si necesitamos imprimir el título de sección
      \ifx\leq@current@section\leq@last@printed@section\else
        \addcontentsline{leq}{leqsection}{\leq@current@section}%
        \global\let\leq@last@printed@section\leq@current@section
        \gdef\leq@last@printed@subsection{}% Reset subsección al cambiar sección
      \fi
      % Verificar si necesitamos imprimir el título de subsección
      \ifx\leq@current@subsection\leq@last@printed@subsection\else
        \ifx\leq@current@subsection\@empty\else
          \addcontentsline{leq}{leqsubsection}{\leq@current@subsection}%
          \global\let\leq@last@printed@subsection\leq@current@subsection
        \fi
      \fi
      % Ecuación
      \refstepcounter{eqnum}%
      \begin{equation*}
        #1\tag{E.\theeqnum}
      \end{equation*}%
      \addcontentsline{leq}{leqt}{[E.\theeqnum] -- #2}%
      \addcontentsline{leq}{leqm}{#1}%
    }
    
    %%% ----- Índice de ecuaciones ----------%%%
    % Tamaño para []
    \newcommand{\leqIndexFont}{\normalsize}
    \newcommand{\leqTitleFont}{\tiny}
    \newcommand{\leqNoteFont}{\tiny}
    % Cabecera de sección 
    \newcommand*\l@leqsection[2]{%
      \addvspace{25pt}% Mayor separación antes de nueva sección
      \noindent\leqIndexFont
      \makebox[\linewidth]{\rule{.38\linewidth}{0.4pt}\ \ \textbf{  \MakeUppercase{#1}}\ \ \rule{.38\linewidth}{0.4pt}}%
      \par\vspace{-10pt}
    }
    % Cabecera de subsección 
    \newcommand*\l@leqsubsection[2]{%
      \par\addvspace{15pt}%
      \noindent\leqIndexFont #1
      \par\vspace{-7pt}
    }
    % Título de ecuación 
    \newcommand*\l@leqt[2]{%
    \par\addvspace{10pt}%
      \par\noindent\leqTitleFont #1\par
    }
    % Miniatura de ecuación
    \newcommand*\l@leqm[2]{%
      \vspace{0.3\baselineskip}%
      \par\scriptsize\noindent\hspace*{1.5em}\(#1\)\par%
    }
    % Nota de ecuación 
    \newcommand*\l@leqnote[2]{%
      \vspace{0.2\baselineskip}%
      \par\noindent\leqNoteFont\hspace*{1.5em}\textbf{Nota.--} #1\par\vspace{0.5\baselineskip}%
    }
    % Generación del índice
    \newcommand{\mylistofequations}{%
      \clearpage
      \phantomsection
      % Título visual de la lista (ajústalo a tu gusto):
      {\centering\bfseries\Large Índice de ecuaciones\par}
      % Entrada en nuestro índice 'stc':
      \addcontentsline{stc}{section}{Índice de ecuaciones}%
      % Llamada a tu lista real (usa el nombre exacto de tu macro):
      \@starttoc{leq}%
    }
    
    %--------- Captura de títulos de sección --------%
    \let\leq@original@section\section
    \let\leq@original@subsection\subsection
    % Flag para saber si estamos en una sección asterisco
    \newif\ifLeqStarSection
    \LeqStarSectionfalse
    
    % Redefinir \section CON CONTROL DE ASTERISCO
    \renewcommand{\section}{%
      \@ifstar{\leq@section@star}{\@dblarg\leq@section@nostar}%
    }
    \def\leq@section@star#1{%
      \global\LeqStarSectiontrue
      \gdef\leq@current@section{#1}%
      \gdef\leq@current@subsection{}%
      \leq@original@section*{#1}%
      \global\LeqStarSectionfalse
    }
    \def\leq@section@nostar[#1]#2{%
      \global\LeqStarSectionfalse
      \gdef\leq@current@section{#2}%
      \gdef\leq@current@subsection{}%
      \leq@original@section[#1]{#2}%
    }
    % Redefinir \subsection
    \renewcommand{\subsection}{%
      \@ifstar{\leq@subsection@star}{\@dblarg\leq@subsection@nostar}%
    }
    \def\leq@subsection@star#1{%
      \global\LeqStarSectiontrue
      \gdef\leq@current@subsection{#1}%
      \leq@original@subsection*{#1}%
      \global\LeqStarSectionfalse
    }
    \def\leq@subsection@nostar[#1]#2{%
      \global\LeqStarSectionfalse
      \gdef\leq@current@subsection{#2}%
      \leq@original@subsection[#1]{#2}%
    }

    % ----- Restauración de valores Globales de estilo ----------%
    % Flags
    \newif\ifInFootnote
    \InFootnotefalse
    \pretocmd{\@footnotetext}{\InFootnotetrue}{}{}
    \apptocmd{\@footnotetext}{\InFootnotefalse}{}{}
    % Profundidad de listas (soporta anidadas)
    \newcount\MyListDepth \MyListDepth=0
    \AtBeginEnvironment{la}{\advance\MyListDepth by 1}
    \AfterEndEnvironment{la}{\advance\MyListDepth by -1}
    \AtBeginEnvironment{lnum}{\advance\MyListDepth by 1}
    \AfterEndEnvironment{lnum}{\advance\MyListDepth by -1}
    \AtBeginEnvironment{itemize}{\advance\MyListDepth by 1}
    \AfterEndEnvironment{itemize}{\advance\MyListDepth by -1}
    \AtBeginEnvironment{enumerate}{\advance\MyListDepth by 1}
    \AfterEndEnvironment{enumerate}{\advance\MyListDepth by -1}
    % Restauración diferida: hazlo al INICIO del próximo párrafo real
    \newif\if@pendingRestore
    \@pendingRestorefalse
    \newcommand{\RequestRestore}{\global\@pendingRestoretrue}
    \everypar{%
      \if@pendingRestore
        \global\@pendingRestorefalse
        \ifInFootnote\else
          \ifnum\MyListDepth=0
            \setlength{\parskip}{\GlobalParskip}%
            \setstretch{\GlobalBaseLineStretch}%
          \fi
        \fi
      \fi
    }
    % Al final de bloques:
    \AfterEndEnvironment{equation}{\RequestRestore}
    \AfterEndEnvironment{equation*}{\RequestRestore}
    \AfterEndEnvironment{la}{\RequestRestore}
    \AfterEndEnvironment{lnum}{\RequestRestore}
    % Al final de macros:
    \apptocmd{\eqblock}{\RequestRestore}{}{}
    \apptocmd{\imagenfit}{\RequestRestore}{}{}
    
\makeatother

% ===================== ToC propio con 4 niveles (único bloque) =====================
\makeatletter

% --- Escritores: añadimos una línea al .stc en cada cabecera numerada ---
% Guardamos las versiones actuales para llamar al comportamiento normal
\let\MyTOC@orig@section\section
\let\MyTOC@orig@subsection\subsection
\let\MyTOC@orig@subsubsection\subsubsection
\let\MyTOC@orig@paragraph\paragraph

% SECTION
\renewcommand{\section}{%
  \@ifstar{\MyTOC@section@star}{\@dblarg\MyTOC@section@nostar}%
}
\def\MyTOC@section@star#1{%
  \MyTOC@orig@section*{#1}% sin entrada al .stc
}
\def\MyTOC@section@nostar[#1]#2{%
  \MyTOC@orig@section[{#1}]{#2}% comportamiento normal (escribe al .toc)
  % Escribimos también nuestra línea al .stc (texto con \numberline)
  \addcontentsline{stc}{section}{\protect\numberline{\thesection}#2}%
}

% SUBSECTION
\renewcommand{\subsection}{%
  \@ifstar{\MyTOC@subsection@star}{\@dblarg\MyTOC@subsection@nostar}%
}
\def\MyTOC@subsection@star#1{%
  \MyTOC@orig@subsection*{#1}%
}
\def\MyTOC@subsection@nostar[#1]#2{%
  \MyTOC@orig@subsection[{#1}]{#2}%
  \addcontentsline{stc}{subsection}{\protect\numberline{\thesubsection}#2}%
}

% SUBSUBSECTION
\renewcommand{\subsubsection}{%
  \@ifstar{\MyTOC@subsubsection@star}{\@dblarg\MyTOC@subsubsection@nostar}%
}
\def\MyTOC@subsubsection@star#1{%
  \MyTOC@orig@subsubsection*{#1}%
}
\def\MyTOC@subsubsection@nostar[#1]#2{%
  \MyTOC@orig@subsubsection[{#1}]{#2}%
  \addcontentsline{stc}{subsubsection}{\protect\numberline{\thesubsubsection}#2}%
}

% PARAGRAPH
\renewcommand{\paragraph}{%
  \@ifstar{\MyTOC@paragraph@star}{\@dblarg\MyTOC@paragraph@nostar}%
}
\def\MyTOC@paragraph@star#1{%
  \MyTOC@orig@paragraph*{#1}%
}
\def\MyTOC@paragraph@nostar[#1]#2{%
  \MyTOC@orig@paragraph[{#1}]{#2}%
  % Si no numeras los \paragraph, cambia \theparagraph por texto plano sin \numberline
  \addcontentsline{stc}{paragraph}{\protect\numberline{\theparagraph}#2}%
}

% --- Impresor del índice propio (lee el .stc) ---
\newcommand{\MyTOC}{%
  \par\bigskip
  {\centering\bfseries\Large Índice de contenido\par}%
  \medskip
  \begingroup
    \parindent=0pt\parskip=2pt
    \@starttoc{stc}%
  \endgroup
}

\makeatother
% ===================== fin ToC propio =====================


%%%%%%%%%%%%%%


% --- Datos del tema ---
\title{Tema 3.A.16 -- Análisis de mercados (I)\\
\large El modelo de competencia perfecta. Análisis de equilibrio parcial: corto y largo plazo; dinámicas de ajuste y estabilidad del equilibrio. Análisis de eficiencia y bienestar.}
\author{Marcelo Ortega}
\date{Febrero 2026}
\newcommand{\TemaCode}{3A16}

\oldtitle{Tema 3.A.14. El modelo de competencia perfecta.}
\changerationale{
    Se desarrolla el título para guiar al opositor sobre el contenido mínimo y garantizar que se cubra en la exposición del tema.
}

\begin{document}


\maketitle
\changebox

\newpage

% --- Página de resumen y modificaciones ---
\puntosclave{ %% Escribir puntos clave Apuntes CECO
    El modelo de mercado en competencia perfecta sirve como marco de referencia ante situaciones más realistas en las que existen fallos de mercado (poder de mercado, externalidades). En todo caso la definición de equilibrio competetivo es análoga a la de cualquier equilibrio económico: todos los agentes toman decisiones óptimas y éstas son compatibles (dada la información disponible para los agentes)

    Permite introducir argumentos gráficos e inequívocos para la medición de las ganancias de bienestar generadas por el intercambio (y compararlas con situaciones con fallos de mercado)
    
    El enfoque de equilibrio parcial permite aislar cualquier fenómeno existente en un mercado de la interrelación con otros mercados. De este manera se permite una interpretación más clara de los efectos de fallos de mercado, y de las medidas que podrían ayudar a solventarlos.
    
    La diferencia entre equilibrio general y equilibrio parcial, es, pues, una cuestión de grado de generalidad en el detalle del análisis, qué grado de generalidad están dispuestos a anunciar como analistas a cambio de entender los mecanismos y los efectos de un fallo de mercado, por ejemplo un enfoque de equilibrio general, nos permiten analizar los efectos de un incremento en el grado de concentración de una industria (p.e. automóviles) en otros mercados (construcción), a través de la relación entre los mercados de inputs. Sin embargo, es muy probable que estos efectos sean de una magnitud mucho menor a los efectos dentro del propio mercado. Por otro lado, cuando la interrelación de los diferentes mercados es alta, el enfoque de equilibrio parcial puede resultar inapropiado, porque los efectos sobre el intercambio y el bienestar pueden afectar a diferentes mercados. Este sería el caso en problemas de introducción de salarios mínimos de deducciones fiscales a rentas, bajas sobre la demanda de vivienda y/o alquiler.
    }
\vspace{1cm}

\modificaciones{
    Última modificación: 11/02/2026

    […] Introducir (al final) como extensión posible:
    \begin{lnum}
        \item Modelo de la Telaraña (equilibrios secuenciales), dinámicas de ajuste al largo plazo; \textcolor{red}{Está en el ANEXO}
        \item \textcolor{red}{Deseconomías de escala externas (por competencia en el mercado de factores productivos)}
    \end{lnum}

    \textcolor{red}{Añadir la parte de largo plazo de los apuntes de Víctor (barreras, estructura de costes, etc.)} \textbf{FORMA PARTE DEL TíTULO DEL TEMA}

    \textcolor{red}{Añadir: TEOREMA DE CHATELIER-SAMUELSON!!!!!!!}
    
    }

\newpage
\MyTOC
\newpage

\mylistofequations
\clearpage

\MyListOfFigures
\clearpage


\pagenumbering{arabic}
\setcounter{page}{1}


\section{Introducción}\label{intro}


\subsection{Contextualización}\label{context}


\subsubsection{Relevancia}\label{importance}


\subsubsection{Historia}\label{history}



\subsubsection{Problemática}\label{problem} 




\subsection{Estructura de la exposición}\label{structure}




\clearpage
\section{El Modelo de Competencia Perfecta}
El punto de partida del análisis es una descripción lo más abstracta, general y sencilla posible de los elementos básicos que constituyen una economía.
\begin{lnum}
    \item Consideramos una economía consistente en unos conjuntos de consumidores ($I$), empresas ($J$), y mercancías ($L$);
    \item Cada consumidor tiene preferencias sobre las cestas de consumo y estas preferencias se pueden representar por una función de utilidad;
    \item Inicialmente, existen unas dotaciones iniciales de las mercancías, pero existe una tecnología disponible para las empresas qué les permite transformar (reduciendo) las dotaciones de algunas mercancías en unidades adicionales de otras, por lo que las cantidades de mercancías finalmente disponibles pueden ser distintas a las dotaciones iniciales. Un plan de producción es factible si puede obtenerse transformando las dotaciones iniciales a través del uso de la tecnología.\footnote{%
        Un plan de producción para cdaa empresa es una lista que especifica cuánto produce (positivo output) o usa (negativo input) de cada mercancia utilizando la tecnología.
    }  
\end{lnum}

Como hemos discutido anteriormente, en este tema nos centraremos en los determinantes de la producción y asignación de un único bien considerado aisladamente, por lo tanto, estaremos tratando con el enfoque analítico parcial introducido por MARSHALL: el equilibrio que se alcanza en un único mercado cometitivo.\footnote{%
    No contemplamos la posibilidad de que se produzcan fenómenos de retroalimentación en las que los cambios en este mercado generen cambios en los mercados de otros bienes, y que terminen afectando al mercado inicial. Esta situación es justificable si el mercado que analizamos es relativamente pequeño con respecto al total de la economía, y no se producen efectos renta en la demanda de otros productos. Así, los cambios en el mercado inicial no provocan cambios en los precios y en las decisiones de los agentes sobre las demás mercancías.
} Por tanto, al mantenerse el precio del resto de bienes fijo, el gasto en el resto de bienes se comporta como un único bien compuesto, al que denominamos numerario:
\begin{la}
    \item \textbf{Preferencias cuasilineales}. Por el lado de los consumidores, la existencia de un bien numerario y la ausencia de efectos renta implica que las preferencias de los consumidores son cuasilineales; y,
    \item \textbf{Tecnología de Producción definida en términos del bien numerario}. Por el lado de las empresas, implica que cada empresa está dotada de una tecnología que le permite transformar el numerario en el bien, de acuerdo a una función de costes que indica el coste mínimo de los inputs necesarios para producir el bien en cuestión dados los precios fijos de los factores. 
\end{la}

Sin pérdida de generalidad, normalizamos el precio del numerario a 1 y denotamos por $p$ el precio del bien no numerario. El intercambio se produce cuando un consumidor está dispuesto a pagar el precio del bien no numerario, que no es más que la cantidad del bien numerario que cualquier empresa acepta a cambio de una unidad del bien no numerario.

En una economía competitiva, consideramos que las empresas y los consumidores individuales son pequeños en relación al tamaño del mercado y sus decisiones individuales no afectan a la determinación del precio del bien no numerario. Sin embargo, el conjunto de las decisiones de consumidores y empresas sí determina el precio de este bien no numerario y la cantidad finalmente intercambiada.

\subsection{La Demanda: el Consumidor}
Consideramos un mercado de un bien, $x\in L$, con un conjunto finito de consumidores $I$ que tienen preferencias cuasilineales sobre combinaciones del bien numerario, $m= (L\setminus\{x\})$, y el bien no numerario $x$. Por lo tanto, las preferencias del consumidor $i$ están representadas por una función de utilidad cuasilineal: 

\eqblock{
\begin{aligned}
    U_i(m,x)=m+\varphi(x)\qquad \text{Por ejemplo:}\quad
    \left\{\begin{aligned}
        U_i(m,x)=m+\ln(x)\qquad (\text{\textcolor{blue}{azul}})\\[0.5em]
        U_i(m,x)=m+\sqrt{x}\qquad (\text{\textcolor{green}{verde}})
    \end{aligned}\right.
\end{aligned}
}{El Problema de Maximización de la Utilidad. Modelo de Competencia Perfecta: Curva de Demanda}
 
Donde, $\varphi(x)$ representa la utilidad en términos numerario  generada por el consumo del bien no numerario.\footnote{%
    “En términos numerarios” significa que la utilidad se mide en las mismas unidades que el bien numerario. En otras palabras: cada unidad de utilidad adicional que proviene del consumo del bien no numerario equivale al valor, en unidades de numerario, que el consumidor estaría dispuesto a pagar (o aceptar) por esa cantidad del bien. Es decir, si $\varphi(x^*)=10\Longrightarrow$ La utilidad proporcionada por el consumo de $x^*$ unidades del bien no numerario (en este caso, 10) equivale a disponer de 10 unidades del bien numerario. Formalmente, dado que la utilidad marginal del numerario es 1 (lo hemos definido así), el numerario sirve como unidad de medida de la utilidad.

Esta función es creciente, cóncava, doblemente diferenciable, y normalizamos de manera que $\varphi(0)=0$.
} Gráficamente, las curvas de indiferencia de los ejemplos dados se podrían representar de la siguiente forma:

\imagenfit{Pref_Cuasilineales_Cons.png}{Preferencias cuasilineales: dos formas funcionales de la Función de Utilidad}{Elaboración propia}

Cada consumidor tiene una dotación inicial del bien numerario, $m_0=W_i$, que puede tener su origen en las rentas obtenidas por la venta de otros recursos (bienes o trabajo dados unos precios fijados y que no se ven afectados por el mercado del bien $x$) o por las rentas derivadas de los beneficios de las empresas que pueda poseer.

\subsubsection{Función de Demanda Individual}
Dada la renta del consumidor, $W_i$, y el precio del bien no numerario, $p^*_x$, la decisión óptima de consumo del bien $x$ del consumidor $i$ se obtendrá como la solución del problema de maximización de utilidad sujeto a la restricción presupuestaria del consumidor por la que el valor del numerario y del bien consumido a los precios de mercado debe ser igual a la renta inicial del consumidor. Se trata de un problema de elección del consumidor clásico (enfoque de maximización de la utilidad).

La condición de optimalidad del consumo de la renta del consumidor, y la demanda óptima del bien sólo depende de su precio.\footnote{%
    Si nos interesara la demanda de varios bienes no numerarios, como en los modelos de diferenciación horizontal o competecia monopolística, la demanda si puede depender del precio de otros productos, pero no se producirán efectos renta.
}


\eqblock{
\begin{aligned}
    &\boxed{
    \begin{aligned}
        &\max_{x}\quad U_i(m,x)\\[0.5em]
        &\text{s.a.}\quad 
        \left\{\begin{aligned}
            &\underset{\text{Restricción Presupuestaria}}{W_i=m^*+p_xx^*}\\[0.5em]
            &\underset{\text{Restricción de No-Negatividad}}{0\ll(m^*,x^*)}
        \end{aligned}\right.
    \end{aligned}
    }\;\underset{
    \text{\scriptsize$
    \left(\begin{aligned}
        &W_i=m_0\;\Rightarrow\;m^*=W_i-p_xx\\[0.5em]
        &\Rightarrow\;m=W_i-p_xx
    \end{aligned}\right)$}
    }{\Longrightarrow}\;
    \boxed{
    \begin{aligned}
        &\max_{x}\quad U_i(x)=W_i-p_xx+\varphi_i(x)\\[0.5em]
        &\text{s.a.}\quad 
        \left\{\begin{aligned}
            &\underset{\text{Restricción de No-Negatividad}}{0\ll x^*}
        \end{aligned}\right.
    \end{aligned}}\\[2em]
    &\Longrightarrow\hspace{1em}\boxed{\frac{\partial}{\partial x}U_i(x)=\frac{\partial}{\partial x}-p_x=\frac{\partial}{\partial x}\varphi(x)=p_x}
\end{aligned}
}{Solución interior al problema de maximización de $U_i$ dada una dotación inicial ($W_i=m_0$). Modelo de Competencia Perfecta: Construcción de la curva de demanda de mercado}

Como vemos, en la solución interior del problema (con consumo de $x^*$ mayor a cero), y sustituyendo el valor del numerario por el obtenido de la restricción presupuestaria, la condición necesaria y suficiente implica que la $UMg$ del bien no numerario sea igual a su precio. 

Esta condición de optimalidad define la curva de demanda inversa del consumidor $i$. Como $\varphi_i'$ es decreciente y positiva, por el Teorema de la Función Implícita podemos obtener la Curva Individual de Demanda Marshalliana (individual), que indica el consumo óptimo para cada precio del bien y que nos permitirá obtener la Función de Demanda Agregada (suma horizontal, por cantidades-$\bar p_x$):

\eqblock{
\begin{aligned}
    &\text{Como:}\quad \frac{\partial}{\partial x}\varphi_i(x):\;\;
    \left\{\begin{aligned}
        &(1)\;\text{Continua}\\[0.5em]
        &(2)\;\text{Decreciente}
    \end{aligned}\right.\qquad\underset{\substack{\text{\scriptsize Teorema de inversión}\\\text{\scriptsize de funciones monótonas}}}{\Longrightarrow}\qquad \exists \left(\frac{\partial}{\partial x}\varphi_i(x)\right)^{(-1)}\equiv g_i(\cdot)\\[2em]
    &\underset{\text{Condición de optimalidad}}{\frac{\partial }{\partial x}\varphi_i(x)=p_x}\qquad \overset{\substack{\text{\scriptsize Aplicamos $g_i(\cdot)$}\\\text{\scriptsize a ambos lados}}}{\Longrightarrow}\qquad g_i\left(\frac{\partial}{\partial x}\varphi_i(x)\right)=g_i(p_x)\\[1em]
    &\text{Obtenemos:}\qquad \underset{\text{Función de Demanda Individual (marshalliana)}}{\boxed{x_i=g_i(p_x)=\left(\frac{\partial}{\partial x}\varphi_i(x)\right)^{(-1)}(p_x)\;\equiv\;x_i(p_x)}}\\[2em]
    &\underset{\substack{\text{\scriptsize Agregación}\\\text{\scriptsize individual}}}{\Longrightarrow}\hspace{4.5em}\underset{\text{\scriptsize Función de Demanda Agregada}}{\boxed{\sum_{i=1}^Ix_i(p_x)=x(p_x)}}
\end{aligned}
}{Obtención e la Función Agregada de Demanda Marchalliana. Competencia Perfecta: Obtención de la Curva de Demanda}
 
Hasta aquí sería el punto compartido con MARSHALL, puesto que se basa en el equilibrio por cantidades. Gráficamente:

\imagenfit{Curva_Demanda_Mercado.png}{Curva de Demanda de Mercado (\textcolor{purple}{morado}): Enfoque parcial}{Elaboración propia}

\subsubsection{Función de Demanda Agregada}
No obstante, la representación (e interpretación) moderna del equilibrio es el ajuste vía precios, por tanto cabe hacer una última operación para obtener la función final.

\eqblock{
\begin{aligned}
    &\text{Como:}\qquad x(p_x):\;
    \begin{aligned}
        \left\{\begin{aligned}
        &(1)\;\text{Continua}\\[0.5em]
        &(2)\;\text{Monótona}
    \end{aligned}\right.\qquad\underset{\substack{\text{\scriptsize Teorema de inversión}\\\text{\scriptsize de funciones monótonas}}}{\Longrightarrow}\qquad \exists \;\left(x(\cdot)\right)^{(-1)}\equiv f(\cdot)\\[2em]
    \end{aligned}\\[2em]
    &\text{Dado $\bar x$}\;:\;x(p_x)=\bar x\qquad\overset{\substack{\text{\scriptsize Aplicamos $f(\cdot)$}\\\text{\scriptsize a ambos lados}}}{\Longrightarrow}\qquad f(x(p_x))=f(\bar x)\qquad \Longrightarrow\qquad \overset{\substack{\sim P_x(\bar x)=x^{(-1)}(\bar x)\\[0.5em]\\}}{\underset{\text{\scriptsize Función de Demanda Inversa}}{\boxed{P^d_x(x)=x^{(-1)}(x)}}}
\end{aligned}
}{Obtención de la Función Inversa de Demanda del Mercado. Competencia Perfecta: Cosntrucción de la Curva de Demanda}

Con esto, obtenemos finalmente la Función Inversa de Demanda de Mercado (\textit{el precio en función de la cantidad demandada}).\footnote{%
    Se introduce la barra sobre \bar x para indicar que estamos evaluando la demanda inversa en un nivel concreto y dado de consumo del bien no numerario, manteniendo fija la renta en bien numerario. Aunque no se escriba explícitamente, la demanda marshalliana (y, por agregación, la demanda de mercado) depende tanto del precio del bien no numerario como de la renta; sin embargo, en este análisis la renta se considera constante. Al fijar la renta, la función de demanda pasa a ser efectivamente univariante en el precio, lo que permite invertirla sin ambigüedad. La barra sirve, por tanto, para dejar claro que la inversión se realiza condicionada a un nivel dado de renta, y que el precio se obtiene como función de una cantidad previamente fijada.
} Asumiendo que todos los consumidores adquieren una cantidad estrictamente positiva del bien no numerario, el beneficio o utilidad marginal para cada consumidor en términos del bien numerario de consumir una unidad adicional del bien no numerario, $\varphi_i'(\uparrow x)$, es exactamente igual a la función inversa de demanda de mercado. Por lo tanto, el valor de la función inversa de demanda del mercado, $P_x(x)$, nos \textbf{indica el beneficio marginal social del consumo del bien no numerario cuando una cantidad agregada del bien es eficientemente distribuida entre los consumidores.}\footnote{%
    Nótese que la eficiencia en la distribución del consumo del bien está asegurada por el hecho de que el beneficio marginal del consumo del bien no numerario es igual para todos los consumidores.

} 

\subsection{La Oferta: la Curva de Oferta del Productor}
El lado de la oferta del mercado está formado por un conjunto de empresas $J$. Cada empresa $j$ está dotado de una tecnología que le permite transformar el numerario en el bien, de acuerdo a una función de costes.\footnote{
En nuestro análisis asumiremos que la función de costes de cada empresa es estríctamente creciente, y doblemente diferenciable para todo nivel de producción estrictamente positivo.  La función de costes de cada empresa puede interpretarse como el resultado de la utilización eficiente de diferentes factores de producción dado un vector de precios fijos de los inputs, de manera que cada empresa minimiza el coste de los factores necesarios para alcanzar cada nivel de producción.

Trabajaremos sobre el supuesto de que esta función de costes es dos veces diferenciables, con primera y segunda derivada mayor que cero. Es decir, la tecnología es estríctamente convexa. El supuesto subyacente es que el mercado de factores productivos es perfectamente competitivo. Sin embargo, como introduciremos más adelante, no tiene porque ser así, en cuyo caso tendremos otros resultados.
} Las empresas son maximizadoras de beneficios y, por tanto, dado un precio, $p^*$, de mercado, el nivel de producción óptimo para una empresa $j$ será el nivel de producción para el que el precio se iguala al coste marginal de producción (en las soluciones interiores), siempre que los beneficios finales sean mayores que el resultado obtenido con un nivel de producción nulo. Formalmente:

\eqblock{
\begin{aligned}
    &\boxed{\begin{aligned}
        &\max_{q_x}\quad \Pi(q_x)=p^*_xq_x-C_j(q_x)\\[0.5em]
        &\text{s.a.}\;
        \left\{\begin{aligned}
            C_j(0)<p^*q_x-C_j(q_x)
        \end{aligned}\right.
    \end{aligned}}\\[2em]
    &\underset{\text{\scriptsize Resolviendo}}{\Longrightarrow}\hspace{2em}
    \left\{\begin{aligned}
        &\text{Condición de máximo:}&&\underset{\text{\scriptsize Debe cumplir esta condición para poder garantizar que estamos ante un máximo}}{\frac{\partial }{\partial q_x}\Pi(q_x)=p^*_x-\frac{\partial}{\partial q_x}C_j(q_x)=0\;\Longrightarrow\;\frac{\partial}{\partial q_x}C_j(q_x)=p^*_x}\\[1em]
        &\text{Condición de restricción:}&&\underset{\text{\scriptsize Debe cumplir esta condición para satisfacer la restricción de participación (o beneficio nulo)}}{\frac{C_j(q_x)-C_j(0)}{q_x}\leq \frac{\partial }{\partial q_x}C_j(q_x)}
    \end{aligned}\right.
\end{aligned}
}{}
 
\subsubsection{Función de Oferta Individual}
De la misma manera que antes, la resolución del problema de maximización planteado nos permite obtener la Función Valor (Función de Producción Indirecta), y con ésta la derivación de la Función inversa de oferta de cada empresa, que representará el precio al que cada empresa estaría dispuesta a ofrecer en el mercado cada nivel de producción. Por tanto, podemos obtener la Función de Oferta de Mercado (Curva de Oferta Agregada) de la siguiente forma:

\eqblock{
\begin{aligned}
    &\text{Como:}\qquad \frac{\partial}{\partial q_x}C_j(q_x):\;
    \left\{\begin{aligned}
        &(1)\;\text{Continua}\\[0.5em]
        &(2)\;\text{Monótona}
    \end{aligned}\right.\qquad\underset{\substack{\text{\scriptsize Teorema de inversión}\\\text{\scriptsize de funciones monótonas}}}{\Longrightarrow}\qquad \exists \;\left(\frac{\partial}{\partial q_x}C_j(q_x)\right)^{(-1)}\equiv g(\cdot)\\[2em]
    &\underset{\text{Condición de optimalidad}}{\frac{\partial }{\partial q_x}C_j(q_x)=p_x^*}\qquad \overset{\substack{\text{\scriptsize Aplicamos $g_i(\cdot)$}\\\text{\scriptsize a ambos lados}}}{\Longrightarrow}\qquad g\left(\frac{\partial }{\partial q_x}C_j(q_x)\right)=g(p^*)\\[1em]
    &\text{Obtenemos:}\qquad q_x=g(p_x^*)\;\Longrightarrow\;\underset{\text{Función de Oferta Individual}}{\boxed{\left(\frac{\partial }{\partial q_x}C_j(q_x)\right)^{(-1)}(p^*)\;\overset{\substack{q_{x,j}\;:\;q_j\\\text{\scriptsize Por simpli-}\\\text{\scriptsize cidad}}}{\equiv}\;q_j(p_x^*)}}
\end{aligned}
}{Obtención e la Función Individual de Oferta. Competencia Perfecta: Obtención de la Curva de Oferta}

Por tanto, para un determinado nivel de producción (dado un determinado nivel de restricción, como antes), la Función de Oferta de la Empresa nos indica cuántas unidades del bien está dispuesta ofrecer la Empresa para un precio determinado.

\subsubsection{Función de Demanda Agregada}
De la misma forma que en el apartado anterior, a patrir de la Función de Oferta de la empresa podemos obtener la Función de Oferta Agregada del Mercado. 

\eqblock{
\begin{aligned}
    &\underset{\substack{\text{\scriptsize Agregación}\\\text{\scriptsize individual}}}{\Longrightarrow}\hspace{4.5em}\underset{\text{\scriptsize Función de Oferta Agregada}}{\sum_{j=1}^Jq_j(p_x^*)=Q(p_x^*)}\\[1em]
    &\text{Como:}\qquad Q(p_x^*):\;
    \begin{aligned}
        \left\{\begin{aligned}
        &(1)\;\text{Continua}\\[0.5em]
        &(2)\;\text{Monótona}
    \end{aligned}\right.\qquad\underset{\substack{\text{\scriptsize Teorema de inversión}\\\text{\scriptsize de funciones monótonas}}}{\Longrightarrow}\qquad \exists \;\left(Q(\cdot)\right)^{(-1)}\equiv f(\cdot)\\[2em]
    \end{aligned}\\[2em]
    &\underset{x\;\sim\; Q_x}{\text{Dado $\bar x$}}\;:\;Q(p_x^*)=\bar x\qquad\overset{\substack{\text{\scriptsize Aplicamos $f(\cdot)$}\\\text{\scriptsize a ambos lados}}}{\Longrightarrow}\qquad f(Q(p_x^*))=f(\bar x)\\[1em]
    &\Longrightarrow\qquad \overset{\substack{\sim P^s_x(\bar x)=Q^{(-1)}(\bar x)\\[0.5em]\\}}{\underset{\text{\scriptsize Función de Oferta Inversa}}{\boxed{P^s_x(x)\equiv Q^{(-1)}(x)}}}\;\sim\;P^s_x(Q)=Q^{(-1)}(Q)
\end{aligned}
}{Obtención de la Función Inversa de Oferta del Mercado. Competencia Perfecta: Cosntrucción de la Curva de Oferta}

Como vemos, en este caso trabajamos también con la Función Inversa de Oferta de Mercado pues expresa el precio al que una empresa estaría dispuesta a vender una cantidad determinada de bienes. Es decir, se realiza un ajuste vía precios y no cantidades. Como cada empresa lleva su producción al punto en que el coste marginal se iguala al precio, el coste marginal de introducir una unidad más de producto en el mercado será igual independientemente de la empresa que lo produzca (y del nivel de producción individual de esta empresa). Por lo tanto, la Función de Oferta Agregada puede considerarse como la Función de $CMg$ de la industria (o la funcón inversa como la función inversa de costes marginales).

\subsubsection{Tecnología de Producción y Curva de Oferta de Mercado}
No obstante, hemos omitido un punto importante en nuestro desarrollo que no podemos seguir sin comentar. ¿Es lógico pensar que todas las funciones de oferta individual son continuas, monótonas y crecientes? Dicho de otra forma, ¿presentan todas las tencologías de producción rendimientos a escala decrecientes? No sería lógica pensarlo y esto es precisamente lo que debemos comentar: existen grandes diferencias a la hora de derivar analíticamente la Función Inversa de Oferta de Mercado en función del tipo de rendimientos que presente la tecnología subyacente.

\paragraph{Rendimientos Decrecientes a Escala}
Si la tecnología de producción presenta una estructura o función de costes estríctamente convexos, la restricción de beneficios mayores que en caso de producción nula se cumple inmediatamente, y:\footnote{%
    Equivalentemente, la tecnología presenta rendimientos marginales decrecientes, de tal forma que el $CMg > CMe$. Podríamos decir que la convexidad de la tecnología de producción en competencia perfecta no es tanto un supuesto, sino una conclusión a la que llegamos; para que se de competenciaperfecta, las empresas deben exhibir rendimientos marginales decrecientes. Por ello hacemos énfasis en que lo importante es que la tecnología presenterendimientos no crecientes a escala, pues rendimientos constantes a escala también pueden presentar rendimientos marginales decrecientes.
} 

\eqblock{
\begin{aligned}
    \text{Si RDE:}\qquad \frac{\partial}{\partial q_x}C_j(q_x)\;\overset{\substack{\text{\scriptsize Función S}\\\text{\scriptsize Individual}}}{\Longrightarrow}\; q_j(p_x^*)\;\overset{\substack{\text{\scriptsize Función S}\\\text{\scriptsize Mercado}}}{\Longrightarrow}\;Q(p_x^*)\;\overset{\substack{\text{\scriptsize Función Inversa}\\\text{\scriptsize Oferta Mercado}}}{\Longrightarrow}\;P_x(x)=Q^{(-1)}(x)
\end{aligned}
}{Tecnología de Producción: Rendimientos Decrecientes a Escala. Modelo de Competencia Perfecta: Obtención Curva de Oferta}

Es, por tanto, equivalente al desarrollo anteriormente realizado. La función de oferta de cada empresa se corresponde con su curva de costes marginales y la condición de maximización de beneficios define inmediatamente la función de oferta de la empresa, que indica el nivel de producción óptimo para la empresa dado un precio. La función de oferta agregada de la industria será, entonces, la suma horizontal de las curvas de oferta de las empresas individuales:

\imagenfit{Ag_S_Mercado.png}{Obtención de la Curva de Oferta de Mercado (FI de Oferta Agregada): Agregación horizontal de las Funciones de Oferta Individuales (RDE)}{\href{https://policonomics.com/es/lp-competencia-perfecta2-oferta-demanda/}{Competencia perfecta II: Oferta y demanda (Policonomics)}}
 
No obstante, tal y como hemos indicado, lo normal en términos analíticos es trabajar con la Función Inversa de Oferta Agregada, puesto que nos permite analizar el ajuste vía precios y no cantidades. En definitiva, nos enconcontraríamos ante el escenario más común en el marco teórico.

\paragraph{Rendimietos Mixtos a Escala: Decrecientes para Niveles Iniciales de Producción}
Trabajar con una estructura de costes decreciente para todos los niveles de producción sería estúpido, puesto que no habría solución al problema de optimización.\footnote{%
    Si la empresa presenta siempre rendimientos crecientes a escala, el problema de maximización de beneficios de la empresa competitiva no tendría solución. Si para un precio $p’$ y la cantidad $q_x’$ la empresa obtuviese beneficios estríctamente positivos, la existencia de rendimientos crecientes a escala implicaría que la empresa siempre obtendrá mayores beneficios aumentando el nivel de producción. Es decir, la existencia de mercados competitivos no es posible en industrias con este tipo de tecnología.
} Por tanto, no analizaremos lo que sucede con una tecnología de prpducción que presenta rendimientos crecientes a escala. No obstante, sí que existe y, de hecho, es común, encontrarnos con estructuras de costes de tipo $U$. Es decir, tecnologías productivas que presentan rendimientos crecientes a escala para niveles iniciales y luego presentan una estructura típica de costes asemejable a las tecnologías de rendimientos decrecientes a escala. 

En presencia de rendimientos crecientes a niveles bajos de producción, la función de costes es inicialmente cóncava, lo que implica costes marginales decrecientes $CMg_j=\frac{\partial}{\partial q_x}C_j(q_x)<0$. En ese tramo, la condición de primer orden $CMg_j(q_x)=p_x^*$ no caracteriza un máximo de beneficios, sino un mínimo, por lo que dichos niveles de producción no son óptimos.\footnote{%
    Como $\Pi_j(q_x)=p_x^*q_x-C_j(q_x)$, si tomamos derivadas e igualamos a cero (CPO), nos queda $\frac{\partial}{\partial q_x}\Pi_j(q_x)=p_x^*-C_j(q_x)$, pero sabemos que, al ser cóncava la función de costes en este punto: $\frac{\partial^2}{\partial^2q_x}C_j(q_x)>0$, lo que significa que cualquier óptimo en este punto representa un mínimo.
} Por ello, una empresa competitiva solo producirá en el tramo en el que la función de costes es convexa, de modo que el coste marginal sea (débilmente) creciente: $CMg_j=\frac{\partial}{\partial q_x}C_j(q_x)\geq0$.

Cuando los costes son cóncavos, se cumple que el coste marginal es menor que el coste medio ($CMg_j<CMe_j$), lo que implica que producir pequeñas cantidades genera pérdidas frente a la alternativa de no producir. Para que los beneficios no sean inferiores a los obtenidos al no producir, es necesario que el precio sea al menos igual al coste medio, lo que solo ocurre a partir del nivel de producción que minimiza el coste medio. Este nivel define la producción mínima eficiente y será el punto a partir del cual se agreguen las Funciones de Oferta Individual para formar la Función de Oferta Agregada:

\eqblock{
\begin{aligned}
    &\text{Si RCrE (iniciales):}\qquad \frac{\partial}{\partial q_x}C_j(q_x)\;\overset{\substack{\text{\scriptsize Dos }\\\text{\scriptsize tramos}}}{\Longrightarrow}\;
    \left\{\begin{aligned}
        &\text{\scriptsize (1) Inferior a escala mínima eficiente}: &&A=\left\{\forall q_x\;:\; \frac{\partial^2}{\partial^2q_x}C_j(q_x)<0\right\}\\[0.5em]
        &\text{\scriptsize (1) Superior a escala mínima eficiente}: &&B=\left\{\forall q_x\;:\; \frac{\partial^2}{\partial^2q_x}C_j(q_x)>0\right\}
    \end{aligned}\right.\\[1em]
    &\boxed{\exists q_j(p^*_x)\iff q_x\in B}\;\Longrightarrow\;
    \left\{\begin{aligned}
        &\forall q_x\in B\;:\;q_j(p_x^*)\;\overset{\substack{\text{\scriptsize Función S}\\\text{\scriptsize Mercado}}}{\Longrightarrow}\;Q(p_x^*)\;\overset{\substack{\text{\scriptsize Función Inversa}\\\text{\scriptsize Oferta Mercado}}}{\Longrightarrow}\;P_x(x)=Q^{(-1)}(x)\\[1em]
        &\forall q_x\not\in B\;\Longrightarrow\;q_j=0,\;Q_x=0
    \end{aligned}\right.
\end{aligned}
}{Tecnología de Producción: Rendimientos Crecientes a niveles iniciales. Modelo de Competencia Perfecta: Obtención Curva de Oferta}

En consecuencia, la curva de oferta individual de la empresa queda determinada del siguiente modo: para cada precio mayor o igual que el mínimo del coste medio, la empresa elige el nivel de producción que iguala el precio al coste marginal en el tramo creciente de este; para precios inferiores, la oferta es nula. Así, la oferta individual coincide con la inversa del coste marginal únicamente en el tramo relevante, esto es, por encima del mínimo del coste medio.

\imagenfit{Cre_NivIniciales_CP.png}{Estructura de Costes de la Empresa: Rendimientos crecientes a niveles iniciales de la producción}{Elaboración propia}

\paragraph{Rendimientos Constantes a Escala}
En el caso de que las funciones de costes de las empresas fuesen lineales y los costes marginales constantes, el problema de maximización sólo admite solución interior si el precio se iguala al coste marginal. Si $p^*_x = \frac{\partial}{\partial q_x}C_j(q_x)$, los beneficios de la empresa serían siempre los mismos para cualquier nivel de producción, y a ese precio cualquier producción maximizaría los beneficios de la empresa.\footnote{%
    Para precios menores al $CMg$, $p^* < CMg$, cualquier incremento de la producción reduciría los beneficios de la empresa y la decisión óptima de la empresa sería una producción nula. Si el precio fuese mayor que el $CMg$, $p > CMg$, la empresa siempre podría aumentar sus beneficios aumentando la producción y por tanto el problema de maximización de los beneficios no admite solución.
} Por lo tanto, la curva de oferta de la empresa sería perfectamente elástica (horizontal) al nivel del coste marginal, y la oferta será nula para precios inferiores al coste marginal. 

\eqblock{
\begin{aligned}
    &\text{Si RCE:}\qquad \frac{\partial}{\partial q_x}C_j(q_x)=\text{cte.}\;\Longrightarrow\;\not\exists\left(\frac{\partial}{\partial q_x}C_j(q_x)\right)^{(-1)}\\[1em]
    &\overset{\substack{\text{\scriptsize Tres}\\\text{\scriptsize posibilidades}}}{\Longrightarrow}\;
    \left\{\begin{aligned}
        &P_x(x)>\text{cte.}\;\Longrightarrow&&\sup_{q_x}\Pi_j(q_x)=\infty\quad\text{(no hay óptimo interior)}\\[0.5em]
        &P_x(x)=\text{cte.}\;\Longrightarrow&&\Pi_j(q_x)=0,\;\;\;\forall q_x\geq0\\[0.5em]
        &P_x(x)<\text{cte.}\;\Longrightarrow&&q_x=0
    \end{aligned}\right.\\[1em]
    &\overset{\text{\scriptsize Mercado}}{\Longrightarrow}\;\;Q(p_x^*):\;
    \left\{\begin{aligned}
        &0&&P_x(x)<\text{cte.}\\[0.5em]
        &\text{cualquier}\;\;\;\underset{Q_x\;=\;x}{Q_x\geq0}&&P_x(x)=\text{cte.}\\[0.5em]
        &\infty&&P_x(x)>\text{cte.}
    \end{aligned}\right.\quad\Longrightarrow\quad \underset{\text{\scriptsize Trabajamos sobre la Función de Oferta \textbf{directa}, no la inversa}}{\boxed{\varepsilon_{Q,P_x}=\frac{\partial Q(p_x)}{\partial p_x}\cdot\frac{p_x}{Q_x}=\infty}}
\end{aligned}
}{Tecnología de Producción: Rendimientos Crecientes a niveles iniciales. Modelo de Competencia Perfecta: Obtención Curva de Oferta}


\clearpage
\section{Equilibrio parcial}
Las curvas de demanda y de oferta (individuales y agregadas) nos indican los planes óptimos de todos los agentes participantes en el mercado cuando este cumple ciertas características deseables: (i) respecto al comportamiento de los agentes; y, (ii) en ausencia de fallos de mercado. Ahora bien, la cuestión que queda por determinar es, precisamemte, el equilibrio en este mercado. Es decir, la combinación precio y cantidad óptima que permite que los planes de los agentes, representados por las funciones agregadas del mercado, coincidan. En este escenario y bajo estos supuestos, el concepto de equilibrio competitivo da respuesta a esta pregunta.

\textbf{Definición.-} Una asignación $X^*=(x_1^*,\dots,x_I^*,q_1^*,\dots,q_J^*)$ y un precio del bien no numerario $p_x^*$ constituyen un equilibrio competitivo si se cumplen las siguientes condiciones:\footnote{%
    Las dos primeras reflejan el supuesto implícito en cualquier modelo económico de que los agentes de la economía buscan tomar las mejores decisiones de acuerdo con sus intereses. Por lo tanto, las empresas maximizan beneficios tomando como dados los precios de los productos y de los inputs, y toman su decisión de acuerdo a su curva de oferta. Los consumidores eligen una cesta de consumo que maximiza su utilidad dada su restricción presupuestar ia, y toman su decisión de acuerdo a su curva de demanda.
}
\begin{lnum}
    \item \textbf{Maximización de Beneficios}. Cada empresa maximiza sus beneficios dados los precios de equilibrio y la tecnología: 
    $$\max \Pi_j(q_x)\iff q_x^*\qquad \text{donde, }\;q_x^*=q_j(p_x^*)$$
    \item \textbf{Maximización de Utilidad}. Cada consumidor maximiza su utilidad sujeto a la restricción presupuestaria:
    $$\max U_i(x)\iff x^*\quad \text{donde,}\;\;x^*=x_i(p_x^*)$$
    \item \textbf{Vaciamiento de Mercados}. La condición de vaciamiento requiere que el precio de equilibrio del plan de consumo y producción sean mutuamente compatibles: la cantidad que los consumidores querrían adquirir a un precio debe igualarse a la producción que las empresas desean poner a la venta a ese mismo precio.\footnote{%
        Si existieran excesos de demanda o de oferta para un precio, la economía no podría estar en equilibrio y por lo tanto, esa situación no podría mantenerse en una economía de mercado. Por ejemplo, si existiese exceso de demanda de un bien, algún consumidor que no puede adquirir una unidad deseada del bien porque no encuentra una empresa que se lo ofrezca a los precios actuales (aunque tiene renta disponible para ello), podría ofrecer un precio ligeramente superior al de mercado y conseguir que algún vendedor ofreciese esa unidad. El argumento simétrico intercambiando los papeles de consumidores y empresa se podría aplicar en el caso de exceso de oferta: Para justificar que en un equilibrio no debe producirse ni exceso de oferta ni de demanda, utilizamos el hecho de que los consumidores y las empresas podrían alterar los precios y, por lo tanto, parece contradecir el supuesto básico de agentes precio aceptantes. Esta aparente paradoja no es tal, ya que en realidad consumidores y empresas siempre tienen esa capacidad (salvo por la existencia de restricciones institucionales como la existencia de precios máximos o mínimos).  La clave del concepto reside en que en el equilibrio los agentes no tendrían incentivos a alterar los precios cuando estos consiguen igualar oferta y demanda. En los precios de equilibrio, los consumidores toman su decisión óptima y no tienen ningún interés en ofrecer mayores precios a las empresas por los productos que adquieren, mientras que las empresas tampoco querrían reducir los precios.
    } Para el precio de equilibrio $p_x^*$ las cantidades totales que los consumidores quieren adquirir del bien no numerario se iguala a las cantidades que las empresas desean producir (equivalentemente si lo plantemos dada una cantidad de bien numerario, $\bar x$):
    $$\begin{aligned}
        &\text{Ajuste cantidades:}&&p_x^*\;:&&x(p_x^*)=Q(p_x^*)\;\iff\;\sum_{i=1}^Ix_i(p_x^*)=\sum_{j=1}^Jq_j(p_x^*)\\[0.5em]
        &\text{Ajuste precios:}&&\bar x\;:&&x^{(-1)}(\bar x)=Q^{(-1)}(\bar x)\;\sim\;P^d_x(\bar x)=P_x^s(\bar x)
    \end{aligned}$$
\end{lnum}

Por tanto, tenemos aquí las tres condiciones que se deben cumplir en el equilibrio de un mercado para definirse como competitivo. 


\subsection{Análisis positivo: Propiedades}
Pasemos a realizar un análisis positivo de este equilibrio, es decir, estudiaremos su propiedades.

\subsubsection{Existencia, Unicidad y Multiplicidad}
Para encontrar el precio de equilibrio del bien, simplemente debemos encontrar el precio para el que la demanda agregada se iguala a la oferta agregada:\footnote{%
    Además, teniendo en cuenta las interpretaciones de las funciones inversas de demanda y de oferta; en el equilibrio competitivo se alcanza un nivel agregado de producción y consumo para el que el beneficio marginal social del consumo del bien no numerario se iguala al coste marginal social de producirlo. Este hecho es un primer indicador de la optimalidad de la asignación del equilibrio competitivo que discutiremos con más detalle más adelante.
}
\begin{la}
    \item \textbf{Curva de demanda}. Bajo los supuestos de continuidad y convexidad estricta de las preferencias de los consumidores aseguradas por la formulación cuasilineal, la curva de demanda será monótona decreciente. Es decir, un aumento del precio conduce a una disminución de la demanda agregada con el fin de mantener la condición de optimalidad: $\frac{\partial}{\partial x}U_i(x)=p_x$. De forma análoga podremos entender lo que sucedería ante una variación en las preferencias de los agentes, o una disminución del precio; y,
    \item \textbf{Curva de oferta}. Por construcción, la curva de oferta existe en el tramo en el cual la tecnología presenta una estructura convexa (estríctamente o no), sabemos que la curva de oferta tendrá pendiente: $\frac{\partial}{\partial Q}C_J(Q_x)\geq 0$
\end{la}

En consecuencia:
\begin{la}
    \item \textbf{Equilibrio único}. El equilibrio puede existir y ser único si:
    $$\exists! \;p_x^*\;:\; x(p_x)=Q(p_x)$$
    En cuyo caso la combinación resultante será de la forma: $\text{Eq}=\{(x_1,\dots,x_I,q_1,\dots,q_J), p_x^*\}$
    \item \textbf{Equilibrio múltiple}. El equilibrio puede existir y ser múltiple si:
    $\exists \:p\;:\;x(p_x)=Q(p_x)$
    En cuyo caso la combinación resultante será de la forma: $\text{Eq}=\{\underset{x_i=0,\;\forall i}{(0,\dots,0},\underset{q_j=0,\;\forall j}{0,\dots,0}),\; \forall p_x^*\in(\max{x^{(-1)}(0)},\dots,\; \min{Q^{(-1)}(0)})\}$
    \item \textbf{Equilibrio inexistente}. El equilibrio no existirá siempre y cuando:
    $\not\exists \;p_x^*\;:\; x(p_x^*)=Q(p_x^*)$
\end{la}
 
\subsubsection{Estabilidad}
La estabilidad del equilibrio parcial competitivo asegura que, a pesar de que no se parta del precio de equilibrio, la economía será capaz de converger a éste. Dos modelos de tiempo continuo se han desarrollado para estudiar la convergencia: el \textit{tâtionnement} y la subasta marshalliana. 

\paragraph{Tâtionnement: Subastador walrasiano}
El modelo de tâtionnement surge con WALRAS. WALRAS asume que los mercados se comportan como si hubiera un subastador que dirigiera el vector de precios hacia los de equilibrio, haciendo uso de la alegoría del proceso de tanteo.\footnote{%
    En realidad, WALRAS utilizó varias veces en su obra \textit{Éléments d’économie politique pure, ou théorie de la richesse sociale} (1874) la palabra tâtonnement (tanteo), ya que el equilibrio se alcanzaría mediante un proceso de ensayo y error, pero en su obra no habla explícitamente del subastador walrasiano, sino que utiliza el pronombre francés \textit{on}. Posteriores autores walrasianos utilizaron la metáfora del subastador walrasiano, pero WALRAS no menciona tal individuo ajeno a la economía.
“ […] supposons qu’on y crie $m\cdot(m−1)$ prix des $m$ marchandises les unes en les autres. Ces prix peuvent être criés au hasard ; pour mieux dire, ils ne peuvent être criés qu’au hasard. Et toutefois il est certain qu’on peut faire en sorte qu’il y en ait parmi eux $(m−1)∙(m−1)=m∙(m−1)−(m−1)$ qui soient liés aux $m−1$ autres par la condition d’équilibre général. A ces prix, ainsi criés, chaque échangeur détermine sa demanda ou offre de (A), (B), (C), (D) [$\cdot$] Cela se fait après réflexion, sans calcul, mais exactement comme cela se ferait par le calcul en vertu du système des équations d’équivalence des quantités demandées et offertes et de satisfaction maximum complété par la restriction convenue. Il s’agit de fonder sur le fait de cette détermination sans calcul une méthode de résolution par tâtonnement des équations d’égalité de l’offre et de la demande totales.”
LEON WALRAS (1874), Éléments d’économie politique pure, ou théorie de la richesse sociale, pág. 127

La idea del tâtonnement no era nueva. De hecho, ya se puede encontrar en La riqueza de las naciones de ADAM SMITH, quien discute la gravitación de los precios de mercado a los precios naturales. Para SMITH, la realización de los precios naturales resulta de la corrección de un estado de desequilibrio en el que los precios efectivos difieren de los precios naturales. Que el comercio tiene lugar a precios de desequilibrio era obvio para SMITH, lo que importaba era la existencia de un mecanismo de reequilibrado. WALRAS pretendía que su noción de tâtonnement recogiera esta idea. Esto implica que el comercio a precios falsos es central en su construcción teórica. 

Desde la perspectiva de WALRAS, el proceso de tâtonnement era una formalización natural de la forma en que cualquier precio competitivo opera en la realidad para establecer un equilibrio competitivo. Tras las críticas de BERTRAND y EDGEWORTH (que afirmaban que al haber intercambio fuera del equilibrio la riqueza de los agentes cambiaba tras cada ronda de negociación), WALRAS admitió en la 4ª edición de sus Elementos que la única salida era basar su análisis en el supuesto de ausencia de comercio fuera del equilibrio, aunque seguía utilizando el término tâtonnement.

La noción de tâtonnement pasó a significar el proceso virtual que ocurre en tiempo lógico (es decir, eliminando todo tipo de comportamiento fuera del equilibrio). Esto implicaba eliminar cualquier tipo de pretensión realista en el análisis del proceso de tâtonnement, que dejo de consistir en un tanteo.
} 

De acuerdo al ajuste walrasiano, la función de exceso de demanda (la diferencia entre las cantidades demandadas y ofertadas al nivel de precios del mercado) depende de los precios.

$$z(p_x)=x(p_x)−Q(p_x)$$

La condición necesaria para la estabilidad es que en caso de excesos de demanda positivos los precios deben subir, y en caso de un exceso de oferta negativo los precios deben bajar.\footnote{
Distinguimos dos tipos de estabilidad: 
\begin{la}
    \item Estabilidad local: Estabilidad en un entorno del equilibrio, por lo que el equilibrio se reestablece tras una perturbación pequeña; y,
    \item Estabilidad global: Estabilidad en todo el conjunto de precios posibles ,por lo que el equilibrio se reestablece ante cualquier perturbación. Esta distinción se realiza por la posibilidad de existencia de equilibrios múltiples.
\end{la}
} 

El subastador canta el precio de un bien en términos de los demás bienes, es decir, el precio del numerario. Si existe un exceso de oferta o un exceso de demanda en alguno de los mercados, no se realiza ninguna transacción, por lo que el subastador debe volver a cantar los precios hasta que se logre el vaciado, momento en el que se da el equilibrio. En caso de exceso de demanda, los precios deben subir y en caso de un exceso de oferta los precios deben bajar. En otras palabras, si dado un precio, la cantidad demandada es mayor que la cantidad ofertada, aumentará el precio. Si por el contrario, dado un precio, la cantidad demandada es menor que la cantidad ofertada, caerá el precio. Es decir, el ajuste se produce vía precios, dando lugar a una convergencia hacia el equilibrio, proceso que le otorga la condición de \textbf{estable}.

\imagenfit{Ajuste_Precios_Walras.png}{Ajuste walrasiano estable}{\href{http://fabiansalazar.es/oposicion/temasenpdf/3A-14.pdf}{Fabián Salazar, M. (2022). Tema 3A-14: El modelo de competencia perfecta.}}

Formalmente, si en $t=0$, $p_x\neq p_x^*$, se requiere que el límite cuanto $t$ tiende a infinito, el precio tienda a aquel vector de precios que lleva al equilibrio, $p_x^*$, es decir: 

$$\lim_{t\to\infty}p_x=p_x^*$$

Por ello, el precio varía con el tiempo, en función del exceso de demanda, siendo $k$ un parámetro de velocidad de ajuste.

\paragraph{Ajuste en cantidades: Subastador marshalliano}
El proceso de Marshall (proceso de subasta de Marshall) adopta un enfoque distinto donde el ajuste se produce \textbf{vía cantidades}. Compradores y productores se encuentran en un mercado donde se inicia una subasta a la holandesa donde los consumidores pujan hasta encontrar el precio que agota las existencia disponibles.

La función de exceso de demanda se define como la diferencia entre los precios de la demanda y oferta en función de la cantidad ofertada en el mercado: 

$$z(x)=P^s_x(\bar x)-P_x^d(\bar x)$$ 

La condición necesaria para la estabilidad es que en caso de $z(x)>0$ (exceso de demanda positivo, entendido como que el precio de demanda es más alto que el precio de oferta) las cantidades deben subir y en caso de $z(x)<0$ las cantidades deben caer. 

En otras palabras, si el precio al que están dispuestos a adquirir el bien los compradores es mayor que el precio al que están dispuestos a vender los productores, aumentará la cantidad intercambiada. Si por el contrario, el precio al que están dispuestos a adquirir el bien los productores es menor que al que están dispuestos a vender los productores, los productores no lo venderán y disminuirá la cantidad intercambiada. 

\imagenfit{Ajuste_Marshall.png}{Ajuste de cantidades hacia el equilibrio en términos diferenciales}{\href{http://fabiansalazar.es/oposicion/temasenpdf/3A-21.pdf}{Fabián Salazar, M. (2022). Tema 3A-21: La teoría del equilibrio general.}}

\subsection{Análisis normativo: Implicaciones de Política Económica}
Habiendo analizado la determinación del equilibrio desde un punto de vista positivo, desde un punto de vista normativo nos podríamos preguntar si la asignación final de recursos constituye un plan socialmente óptimo de consumo y producción en un mercado concreto. Es decir, podríamos valorar la deseabilidad de ese equilibrio, campo que corresponde a la Economía del Bienestar.

\subsubsection{Eficiencia: Asignación de recursos}
El primer aspecto que debemos valorar es la eficiencia. Pero, ¿qué es la eficiencia? El término, es desde luego, muy general, pero podemos definirlo como la \textbf{capacidad de lograr los resultados deseados con el mínimo posible de recursos}. De esta definición extraemos la primera conclusión: la eficiencia y su importancia se desprende de uno de los supuestos económicos más importantes, la escasez.\footnote{%
    Hablaremos en mayor medida de la \textbf{eficiencia} en el [\authorfont{Ver Tema 3.A.22}], pero podemos hacer una introducción breve aquí.

Por otro lado, debemos tener claro que la escasez es el pilar fundamental del verdadero problema económico: \textbf{la elección}. Sin escasez no habría elecciones/decisiones mejores que otras bajo el mismo marco exiomático, puesto que no habría conjunto alguno que restringiese nuestras posibilidades de elección.
} Por tanto, ¿cómo podemos valorar la eficiencia del equilibrio competitivo? ¿qué restricción de recursos se impone? Para resolver estas preguntas debemos acordarnos de uno de los supuestos de partida: la dotación de bien numerario.\footnote{%
    En este enfoque de equilibrio parcial, el bien numerario aglutina, como hemos dicho, no sólo todo el resto de bienes disponibles en la economía (bien sea de consumo o insumos), sino también \textit{bienes} (en abstracto) que propocionan utilidad al consumidor \textit{per se}, como por ejemplo, el ocio, educación, etc. Además, el bien numerario es también el insumo productivo que utilizan las empresas para ejecutar su actividad.
} Por tanto, he aquí la escasez a partir de la cual evaluar la eficiencia asignativa del equilibrio competitivo.

Por tanto, dado que el bien numerario es el recurso básico de la economía, dotación que determina y restringe en última instancia las capacidades productivas de ésta, diremos que el equilibrio es eficiente si y solo si la última unidad del bien numerario consumida para producir el bien no numerario genera el mismo valor que si permaneciera inutilizado. En otras palabras, la eficiencia asignativa se alcanza cuando el numerario se emplea en producción hasta el punto en que el valor marginal de transformar numerario en bienes coincide con el valor marginal de mantenerlo en consumo:

\eqblock{
\begin{aligned}
    &\text{Sea,}\;\;m=\sum_{i=1}^IW_i &&\equiv\text{Restricción de recursos global}\\
    &\text{Consumo:} &&P_x^d(x)\equiv{BMgS}_x\qquad  x^*\Longrightarrow\left.\frac{\frac{\partial}{\partial x}U(x,m)}{\frac{\partial}{\partial m}U(x,m)}\right|_{x=x^*}\overset{(\partial U/\partial m=1)}{=}\frac{P_x(x^*)}{1}=P_x(x^*)\\[0.5em]
    &\text{Producción:} &&P_x^s(x)\equiv{CMgS}_x\qquad  x^*\Longrightarrow\underset{\substack{\text{\scriptsize El coste marginal ya es la productividad}\\\text{\scriptsize del numerario en producción}}}{\left.\frac{\partial}{\partial x}C(x)\right|_{x=x^*}}\\[2em]
    &\underset{\substack{\text{\scriptsize Eficiencia}\\\text{\scriptsize asignativa}}}{\Longrightarrow} &&{BMgS}_x={CMgS}_x\iff \left.\frac{\frac{\partial}{\partial x}U(x,m)}{\frac{\partial}{\partial m}U(x,m)}\right|_{x=x^*}=\left.\frac{\partial}{\partial x}C(x)\right|_{x=x^*}
\end{aligned}
}{Eficiencia asignativa: Restricción impuesta por el numerario. Análisis normativo: Eficiencia y Bienestar}

Por tanto, en un mercado competitivo con bien numerario como recurso productivo, la eficiencia asignativa implica que el numerario se distribuye entre consumo y producción de modo que el valor marginal para los consumidores de la última unidad producida iguala el coste marginal de transformarla, es decir, $P_x(x^*)=CMg(x^*)$.

\paragraph{Excedente Total: MARSHALL}
De forma análoga, una alternativa propuesta por MARSHALL haciendo uso del concepto \textbf{excedente} nos llevará a la misma conclusión: el equilibrio competitivo constituye una asignación eficiente de recursos y \textit{beneficios}.\footnote{%
    Marshall quiere medir ganancias económicas generadas por el intercambio en un mercado, sin entrar en comparaciones interpersonales de utilidad (que él considera problemáticas).
Para ello necesita una magnitud monetaria observable. Si una unidad se produce cuando el beneficio marginal social (disposición a pagar) es mayor o igual que el coste marginal social, esa unidad aumenta el excedente total. El equilibrio competitivo produce exactamente todas esas unidades y ninguna más, por lo que maximiza el excedente agregado. 

Esto no es un criterio paretiano puro, porque el criterio paretiano no permite comparar asignaciones cuando hay ganadores y perdedores. Aquí sí se comparan, implícitamente aceptando que las ganancias de unos pueden compensar las pérdidas de otros. Por eso el criterio subyacente es cardinal y utilitarista, aunque Marshall lo presente de forma pragmática.
} 

En el entorno de equilibrio parcial, definimos el excedente total como las ganancias netas generadas por la producción y el consumo de $x^*$ unidades del bien no numerario. Por tanto, para comprobar si el equilibrio competitivo alcanza el mayor nivel de bienestar calcularemos si, en este, el Excedente Total se maximiza. Este excedente está compuesto por:
\begin{la}
    \item \textbf{Excedente del consumidor}.\footnote{%
        Este concepto fue formalizado por MARSHALL (1890), sin embargo, el concepto ya había sido propuesto anteriormente por el autor proto-marginalista JULES DUPUIT en su obra \textit{De la Mesure de l´Utilité des Travaux Publics} (1884) [\authorfont{Ver Tema 3.A.9}].
    } El excedente del consumidor se define como la diferencia entre lo que el consumidor está dispuesto a pagar por una unidad del bien y lo que finalmente paga; y,
    \item \textbf{Excedente del productor}. El excedente del productor es la diferencia entre el precio al que está dispuesto a vender el productor cada unidad y el precio de mercado al que realmente las vende. Indica la ganancia neta de bienestar obtenida por los productores al producir q unidades del bien no numerario. Esta mide la suma de los beneficios operativos que se generan en el mercado sin tener en cuenta los costes fijos.
\end{la}

Por lo tanto, de forma conjunto podemos expresar analíticamente el método como:

\eqblock{
\begin{aligned}
    &\text{ET}_x=\text{EC}_x+\text{EP}_x\;\Longrightarrow\;\text{ET}_x=\underbrace{\left[\int_{0}^{x}P^d_x(x)\mathrm dx-p_x\cdot x\right]}_{\text{\scriptsize Excedente del Consumidor}}+\underbrace{\left[p_x\cdot x-\int_0^{x}P_x^s(x)\mathrm dx\right]}_{\text{\scriptsize Excedente del Productor}}\\[2em]
    &\text{Bienestar}\;\Longrightarrow\;\boxed{\max \text{ET}_x=\int_0^{x}(P^d_x(x)-P^s_x(x))\mathrm dx}
\end{aligned}
}{Análisis de Bienestar: Excedentes. Analísis normativo: Bienestar}

Tanto gráfica como analíticamente resulta evidente que el valor del excedente total se maximiza en el nivel de consumo $x^*$ para el que $P_x^d(x^*)=P^s_x(x^*)$. Este hecho es una primera constatación del Primer Teorema Fundamental de la Economía del Bienestar (1TFEB), que establece que la asignación del equilibrio competitivo es óptima en el sentido de Pareto y que discutimos en profundidad en [\authorfont{Ver Tema 3.A.22}].

\imagenfit{ET_Eq.png}{Maximización del Excedente Total}{\href{http://fabiansalazar.es/oposicion/temasenpdf/3A-14.pdf}{Fabián Salazar, M. (2022). Tema 3A-14: El modelo de competencia perfecta.}}

\subsubsection{Implicaciones de Política Económica: Bienestar}
Hasta aquí hemos visto como el equilibrio competitivo constituye una asignación eficiente de recursos productivos y, de manera más general, un óptimo de PARETO (ya que no se puede mejorar la situación de uno sin empeorar la de otro agente). 

No obstante, este criterio no resulta suficiente para valorar la deseabilidad del equilibrio, ¿no? Es decir, hemos visto que sí cumple con eficiencia técnica y asignativa, ¿pero es esto suficiente?\footnote{%
    Perfectamente podemos estar ante un escenario en el que una asignación eficiente no es óptimo de PARETO. Por ejemplo, piensese en un monopolio discriminador de primer grado. En el equilibrio bajo estas circunstancias, el equilibrio final cumplirá con la condición de eficiencia asignativa. No obstante, el productor acaparará todo el excedente del consumidor.
} 

Para el planificador público, un criterio de valoración fundamental es el de equidad., aunque en términos económicos estrictos no está definido. Sin embargo, si podemos introducir el criterio de equidad en la distribución final teniendo en cuenta un elemento fundamental: la dotación inicial de factores. Como hemos demostrado antes, el equilibrio competitivo alcanza la eficiencia en el sentido de PARETO, independientemente la distribución dotacional inicial. Ahora bien, por esta misma razón, resulta lógico pensar que el criterio de equidad puede introducirse mediante una redistribución inicial de las dotaciones iniciales. 

\paragraph{2TFEB: Redistribución de las dotaciones iniciales ${\to}$ Óptimo de PARETO}
Dado un vector de utilidades $\vec u=(\bar u_i)_{i\neq1}$, podremos plantear el problema de maximización del planificador benevolente con el que caracterizaremos el conjunto de asignaciones eficientes en el sentido de PARETO, en función de unos multiplicadores (los Lagangianos) que contendrán la información marginal relativa a la escasez.\footnote{%
    En el [{\authorfont{Anexo \ref{anx:primal_planner}}}] podemos ver el desarrollo completo de dicho problema y los principales resultados.
} Una vez que identifiquemos la asignación eficiente que a su vez cumple con las expectativas distributivas del planificador, la pregunta será: ¿existe un sistema de precios y transferencias tal que, si cada agente maximiza individualmente, la asignación escogida por el planificador emerge como equilibrio competitivo?

La respuesta es que sí. Atendiendo al 2TFEB, sabemos que cualquier asignación eficiente en el sentido de PARETO se puede alcanzar mediante un sistema predefinido de precios y transferencias (lo que es equivalente a decir: una redistribución inicial de las dotaciones de recursos).

Por tanto, las implicaciones de Política Económica son claras y abrumadoras. El equilibrio competitivo conduce a una asignación eficiente de los recursos escasos de la economía y, además, mediante una estructura competitiva de mercado podremos alcanzar una asignación eficiente en el sentido de PARETO que atienda a las necesidades distributivas que se definan previamente.

Por lo tanto, transferencias adecuadas de las dotaciones iniciales del bien numerario permitirían alcanzar cualquier asignación factible que maximiza la suma de las utilidades de los consumidores como resultado del equilibrio competitivo. Una autoridad central benevolente interesada en alcanzar una asignación óptima de Pareto particular puede implementar esa asignación deseable mediante transferencias previas del bien numerario entre los consumidores y dejando trabajar al mercado. 

Los mercados perfectamente competitivos son, pues, el estándar sobre el que se deberían comparar los resultados de economías de mercado no perfectamente competitivas. En particular, cualquier ineficiencia que puede surgir en una economía de mercado, y por tanto el papel de intervenciones que supongan mejoras en el sentido de Pareto, deben tener su origen en violaciones de los supuestos explícitos del paradigma competitivo.

\subsection{Estática comparativa}
Una vez hemos definido un modelo que ofrece un concepto positivo que explica un fenómeno partiendo de unos condiciones iniciales, resulta de interés conocer cómo cambiarían las predicciones del modelo ante cambios de los parámetros iniciales. 



Analizaremos brevemente como afectan al equilibrio competitivo cambios en las condiciones iniciales de mercado (preferencias, costes). De esta manera, compararemos los resultados en mercados similares pero no exactamente iguales o conocer los efectos de cambios en las condiciones de mercado en los resultados del equilibrio competitivo.
En general, las alteraciones de mercado que, \textit{a priori}, no pueden catalogarse como fallor de mercado son cambios que afectan a:
\begin{la}
    \item \textbf{Cambio en preferencias} (curva de demanda). Supongamos que, para cada consumidor la utilidad generada por el bien no numerario depende de un parámetro $\alpha$, $\varphi_i(x,\alpha)$. Por ejemplo, cuanto mayor sea la temperatura media de una ciudad, mayor será la utilidad que genera en cada individuo el consumo de helados (el efecto puede ser positivo o negativo);\footnote{%
        No consideraremos cambios en los precios de los demás productivos porque el precio del numerario lo mantenemos constante en uno. Por otro lado, cambios en el nivel de renta tampoco afectan a la demanda del bien no numerario y por lo tanto no generan shocks de demanda.
    }
    \item \textbf{Cambio en tecnología y funciones de costes} (curva de oferta). Suponemos que la función de costes de cada empresa depende de un parámetro $\beta$, $C_j(q_x,\beta)$. Por ejemplo, cambios técnicos o aparición de nuevos recursos pueden reducir los costes de producción;
    \item \textbf{Existencia de impuestos y subsidios}. La existencia de estos elementos implica que los precios efectivos para consumidores y empresas difieren. Así, denotamos por $\widehat{P_x^s}(x)$ y $\widehat{P_x^d}(x)$ los precios efectivos para consumidores y empresas respectivamente, dado el impuesto o subsidio $t$.
\end{la}

\subsubsection{Curva de Demanda}
La curva de demanda se desplazará por cambios en las preferencias o impuestos/subvenciones.

\paragraph{Cambios en las preferencias: $\alpha$}
Desde un punto de vista positivo, un shock positivo de demanda implicará un desplazamiento hacia arriba de la curva de demanda, provocando un aumento de los precios y de la cantidad intercambiada.

\imagenfit{Cambios_Preferencias.png}{Shock positivo de demanda}{Apuntes ICEX-CECO. Tema 3.A.16}

A nivel normativo, esto implicará que ante shocks positivos de demanda, aumentan los precios y la cantidad consumida. Esto se traduce en que el excedente del productor aumentará y el excedente del consumidor quedará indeterminado. En todo caso, el excedente total aumentará.

\paragraph{Impuesto sobre el consumo}
Consideremos la introducción de un impuesto sobre las ventas de manera que los consumidores deben pagar un impuesto unitario $t>0$ por cada unidad del bien de interés adquirida. 

La introducción de este impuesto no afecta ni a las funciones de utilidad de los consumidores ni a las funciones de costes de las empresas, pero sí afectaría a los precios efectivos que pagan los consumidores y que difieren del que reciben las empresas. Así, los consumidores pagan un precio por unidad de producto mayor que el que reciben las empresas, y la diferencia es exactamente la cuantía del impuesto unitario $t$. 

En esta situación, los consumidores perciben un precio $\widehat{P_x^d}(x)=P^d_x(x)+t$, mientras que las empresas reciben $\widehat{P_x^s}(x)=P^s_x(x)$. Las condiciones de optimalidad de los consumidores con la introducción de este impuesto obligan que la demanda agregada para cualquier precio sea necesariamente $x(p_x+t)$, ya que para los consumidores la introducción del impuesto implica que el precio se incrementa en la cuantía del impuesto. Por lo tanto, el precio de equilibrio será, por la condición de vaciamiento de mercado, el que satisfaga:

$$
\begin{aligned}
    &\underset{\text{\scriptsize Impuesto \textit{ad-valorem}, no proporcional}}{\text{Impuesto sobre el consumo:}}&&x(p_x+t)=Q(p_x)
\end{aligned}
$$

\imagenfit{EC_Impuesto.png}{Introducción de un impuesto fijo sobre el consumo unitario}{Apuntes ICEX-CECO. Tema 3.A.16}

Para determinar los efectos en los precios efectivos provocados por un cambio marginal en el impuesto:

\eqblock{
\begin{aligned}
    &\underset{\text{\scriptsize Impuesto \textit{ad-valorem}, no proporcional}}{\text{Impuesto sobre el consumo:}}\qquad \underset{\text{\scriptsize Equilibrio del mercado}}{x(p_x+t)=Q(p_x)}\\[2em]
    &\text{Entonces, podemos escribir:}\qquad  p_x^*\equiv p_x(t)\\[1em]
    &\Longrightarrow\;\text{si,}\qquad \exists\left\{\left(\frac{\partial}{\partial p_x}x(\cdot)\right),\;\;\left(\frac{\partial}{\partial p_x}Q(\cdot)\right)\right\}\\[1em]
    &\Longrightarrow\quad
    \left\{\begin{aligned}
        &\frac{\partial}{\partial t}x(p^*_x(t)+t)=x(p_t^*(t)+t)\cdot \frac{\partial}{\partial t }(p_x^*(t)+t)\Longrightarrow\frac{\partial}{\partial t}(p_x^*(t)+t)=1\\[0.5em]
        &\frac{\partial}{\partial t}Q(p_x^*(t))t=Q(p_x^*(t))\cdot p_x^*(t)
    \end{aligned}\right.\\[0.5em]
    &\text{Impacto:}\qquad x(p^*_x(t)+t)\cdot[p^*_x(t)+1]=q(p_x^*(t))p_x^*(t)\quad\Longrightarrow\quad \boxed{p_x^*=\frac{x(p_x^*(t)+t)}{x(p_x^*(t)+t)-q(p_x^*(t))}}
\end{aligned}
}{Equilibrio tras la introducción de un impuesto sobre el consumo. Estática comparativa: Curva de Demanda}

Desde un punto de vista positivo, cómo y cuánto varían los precios a los que se enfrenta cada uno de los agentes depende de las elasticidades de la oferta y la demanda:\footnote{%
    Cuanto mayor sea la elasticidad de la oferta (más horizontal) menor será la caída en el precio recibida por las empresas y por tanto casi todo el efecto del impuesto recaería sobre los consumidores. Podríamos afirmar que el impuesto recae más sobre la parte más inelástica del mercado.
} (i) el precio efectivo recibido por las empresas cae (débilmente); (ii) el precio efectivo pagado por los consumidores sube (débilmente); y, (iii) la cantidad total intercambiada cae (débilmente).

Desde un punto de vista normativo, buscamos estudiar cómo y cuánto afecta al bienestar. La introducción del impuesto genera una distorsión del sistema de precios en el mercado, ya que la valoración marginal social y el coste marginal social ya no coinciden:
\begin{la}
    \item El precio efectivo recibido por las empresas es menor que cuando no había impuesto, pero la reducción es menor que t; y,
    \item El precio efectivo pagado por los consumidores sube, pero la subida es menor que t.
\end{la}

Por lo tanto, se produce una traslación del impuesto: aunque el impuesto sobre el consumo, su impacto se distribuye entre ambos lados del mercado (ambos pagan parte del impuesto). 

\imagenfit{Impuesto_EC_Consmp.png}{Introducción de un impuesto sobre el consumo unitario}{Apuntes ICEX-CECO. Tema 3.A.16}

Derivando se puede hallar que proporción del impuesto recaerá sobre los consumidores y que parte recaerá sobre los productores:

\eqblock{
\begin{aligned}
    \text{\% Consumidores}:&&\frac{\partial}{\partial t}(p_x^*(t)+t)=\frac{\varepsilon_{S,p_x}}{\varepsilon_{S,p_x}+|\varepsilon_{D,p_x}|}\cdot 100\\[2em]
    \text{\% Productores}:&&\frac{\partial}{\partial t}p_x^*(t)=\frac{|\varepsilon_{D,p_x}|}{\varepsilon_{S,p_x}+|\varepsilon_{D,p_x}|}\cdot 100
\end{aligned}
}{Incidencia del Impuesto sobre el consumo unitario del bien. Estática comparativa: Curva de Demanda}

Esto es lo que se conoce como la ley de DALTON: la incidencia económica de un impuesto en un mercado competitivo se reparte entre oferentes y demandantes de acuerdo a la elasticidad-precio de la oferta y la demanda, con independencia de sobre quién recaiga la incidencia legal, de manera que el impuesto recaerá más sobre el agente cuya elasticidad-precio sea relativamente menor en valor absoluto.

Todo lo anterior genera una disminución tanto en el excedente del consumidor como en el excedente del productor, de modo que se reduce el excedente total. Esa disminución se conoce como la pérdida irrecuperable de eficiencia generada por la distorsión impositiva.\footnote{%
    El modelo cuasilineal también ofrece una justificación de la utilización del excedente total como medida del bienestar generado por potenciales mejoras en el sentido de Pareto. Volviendo al ejemplo de los esquemas de impuestos y transferencias, un cambio en esta política implicaría una mejora potencial si todos los consumidores mejorasen tras la implantación de las transferencias. En el modelo cuasilineal esto sería posible si el cambio en el excedente total fuese positivo. Este cambio positivo implica que las ganancias obtenidas por los que mejoran con la nueva medida de política económica es suficiente para compensar a los perjudicados. Este razonamiento que sirve de justificación para el análisis coste beneficio es conocido como el Principio de Compensación.

El Principio de Compensación permite comparar la optimalidad de diferentes medidas económicas que podrían alterar el conjunto de asignaciones factibles (Frontera de Posibilidades de Bienestar) y no sólo la distribución del bienestar entre los miembros de la sociedad. En los casos generales en los que la comparabilidad de la utilidad entre agentes no es posible, los criterios de compensación no pueden ofrecer criterios completos para comparar diferentes Fronteras de Posibilidades de Bienestar. Sin embargo, en el modelo cuasilineal, cambios en la Frontera de Posibilidades de Bienestar no alteran la relación de intercambio entre las utilidades de los consumidores (1 a 1). Por tanto, el Principio de compensación equivale a la maximización de la suma de utilidad de los consumidores, es decir, el excedente total.
}











\clearpage
\section{Conclusión} 

\subsection{Extensiones y relación con otras partes del Temario}


\subsection{Opinión}



\subsubsection{Idea final (Salida o cierra)}


\clearpage
\pagenumbering{gobble}
\MyBibliography
\MyBibItem{Autor}{Título}{Año}{DOI -- LINK}




%ANEXOS
\clearpage
\anexo{\textbf{Preguntas Test}}
%\input{\TemaCode/questions.tex} 





\clearpage
\anexo{\textit{Primal Planner} -- Análisis del Bienestar (2TFEB): Enfoque parcial}\label{anx:primal_planner}
El problema del planificador se usa para caracterizar el conjunto de asignaciones factibles y Pareto-eficientes a partir de condiciones necesarias de optimalidad (CPO's). Al maximizar una agregación de utilidades sujeta a las restricciones de recursos (y a la factibilidad tecnológica), el Lagrangiano introduce multiplicadores asociados a cada restricción. Esos multiplicadores son precios sombra: miden el valor marginal (en términos del objetivo del planner) de relajar una restricción de recursos. En equilibrio interior, las KKT implican igualación entre tasas marginales de sustitución/valoración social y esos precios sombra, y entre costes marginales y esos mismos precios sombra; esto produce una condición de eficiencia del tipo “beneficio marginal social = coste marginal social” (en el bien no numerario) y, si hay varios bienes/factores, el conjunto completo de condiciones de eficiencia.

$$\begin{aligned}
    &\max_{\{m_i,x_i,q_j,z_j\}}U_1(m_1,x)=m_1+\varphi_1(x_1)\\[0.5em]
    &\text{s.a.}\;
    \left\{\begin{aligned}
        &\underset{\text{\scriptsize Restricción no empeoramiento $\forall i$}}{\bar u_i\leq m_i+\varphi_i(x_i)},&&\forall i\neq1\\
        &\underset{\substack{\text{\scriptsize Restricción de producción}}}{\sum_{i=1}^Ix_i\leq\sum_{j=1}^Jq_j}\\
        &\underset{\text{\scriptsize Restricción de factabilidad: escasez}}{\sum_{i=1}^Im_i+\sum_{j=1}^J\leq\sum_{i=1}^J\omega_i}\\
        &\underset{\text{\scriptsize Restricción tecnológica: no gratis}}{C_j(q_j)\leq z_j,}&&\forall j
    \end{aligned}\right.\\[1em]
    &\begin{aligned}
        &\text{Planteamos Lagrangiano:}\\[0.5em]
        &\mathcal L=m_1+\varphi_1(x_1)+\underset{\text{\scriptsize mínimo garantizado}}{\sum_{i\neq1}\lambda_i(U_i-\bar u_i)}+\underset{\text{\scriptsize vaciamiento de mercado}}{\mu\left(\sum_{j}q_j-\sum_{i}x_i\right)}+\underset{\text{\scriptsize restricción dotacional real}}{\nu\left(\sum_i\omega_i-\sum_im_i-\sum_jz_j\right)}+\underset{\text{\scriptsize no producción gratuita}}{\sum_{j}\eta_j(z_j-C_j(q_j))}\\[1em]
        &\text{Soluciones (CPO's):}
        \left\{\begin{aligned}
            &(1: \partial m_i)\;:\;
            \left\{\begin{aligned}
                &\frac{\partial}{\partial m_1}\mathcal L=1-\nu=0\Longrightarrow\;\nu=1\\
                &{\frac{\partial}{\partial m_i}\mathcal L}_{i\neq1}=\lambda_i-\nu=0\Longrightarrow\;\lambda_i=1
            \end{aligned}\right.\\[0.5em]
            &(2: \partial x_i)\;:\;
            \left\{\begin{aligned}
                &\frac{\partial}{\partial x_1}\mathcal L=\frac{\partial}{\partial x_1}\varphi_1(x_1)-\mu=0&&\Longrightarrow\;&&\frac{\partial}{\partial x_1}\varphi_1(x_1)=\mu\\
                &{\frac{\partial}{\partial x_i}\mathcal L}_{i\neq1}=\lambda_i\left(\frac{\partial}{\partial x_i}\varphi_i(x_i)\right)-\mu=0&&\underset{\text{\scriptsize pero, $\lambda_i=1$}}{\Longrightarrow}\;&&\frac{\partial}{\partial x_i}\varphi_i(x_i)=\mu
            \end{aligned}\right.\\[0.5em]
            &(3: \partial z_j)\;:\;
            \left\{\begin{aligned}
                &{\frac{\partial}{\partial z_j}\mathcal L}_{i\neq1}=\eta_j-\nu=0&&{\Longrightarrow}\;&&\eta_j=\nu=1
            \end{aligned}\right.\\[0.5em]
            &(4: \partial q_j)\;:\;
            \left\{\begin{aligned}
                &{\frac{\partial}{\partial q_j}\mathcal L}_{i\neq1}=\mu-\eta_j\left(\frac{\partial}{\partial q_j}C_j(q_j)\right)=0&&\underset{\text{\scriptsize pero, $\eta_j=1$}}{\Longrightarrow}\;&&\mu=\frac{\partial}{\partial q_j}C_j(q_j)
            \end{aligned}\right.
        \end{aligned}\right.
    \end{aligned}\\[2em]
    &\underset{\text{\scriptsize Resultado}}{\Longrightarrow}
    \left\{\begin{aligned}
        &(m^*,x^*, q^*)\\[0.5em]
        &\text{Definiendo:}\;p^*\equiv\mu\Longrightarrow
        \left\{\begin{aligned}
            &\frac{\partial}{\partial x_i}\varphi_i(x_i^*)=p^*,\;\;\forall i\;\text{con}\;\;x_i^*>0\\[0.5em]
            &\frac{\partial}{\partial q_j}C_j(q_j^*)=p^*,\;\;\forall j\;\text{con}\;\;q_j^*>0
        \end{aligned}\right.\;\;\Longrightarrow\;\boxed{\sum_{i}x_i^*=\sum_jq_j^*}
    \end{aligned}\right.
\end{aligned}$$

Resulta interesante comentar la relación entre el multiplicador de Lagrange $\mu$ y el precio del equilbrio competitivo $p^*$ en las condiciones de maximización del excedente total y las condiciones de optimalidad del comportamiento de consumidores y empresas. El precio del equilibrio competitivo se iguala al precio sombra de la restricción de recursos para el bien no numerario. En ese sentido, el precio del equilibrio competitivo refleja el valor social marginal del bien no numerario (beneficio marginal social). Cada empresa, al llevar el coste marginal de su producción hasta el punto en que se iguala al precio, iguala su coste marginal de producción con el valor marginal para la sociedad de esa producción. Esta correspondencia entre precios de equilibrio y precios-sombra óptimos se verifica generalmente y no es consecuencia de la especificación cuasilineal.

Proposición.- Para cualquier vector de utilidades de los consumidores correspondiente con una asignación óptima de Pareto, existen transferencias del numerario entre los consumidores tal que en el equilibrio competitivo alcanzado en la economía, las utilidades de los consumidores coinciden con las del vector de utilidades óptimo de Pareto inicial.

Esto se puede demostrar de la siguiente forma. Sea una asignación eficiente concreta, esta puede ser realizada por una economía descentralizada, fijando un precio ($p^*$) y se pide que: (i) cada consumidor maximice su utilidad sujeto a su restricción presupuestaria a esos precios; (ii) cada empresa maximice beneficios con su tecnología; y, (iii) los mercados vacíen. El vínculo matemático es que, si tomas p igual al vector de multiplicadores del planner (normalizado por el numerario), entonces las condiciones de primer orden de consumidores y empresas coinciden con las CPO's del planificador para esa asignación: las decisiones individuales generan exactamente los mismos $x_i^* y q_j^*$.

$$\begin{aligned}
    
\end{aligned}$$

La razón por la que aparece un “paso con transferencias” es estrictamente institucional: el planner impone factibilidad agregada, pero no respeta presupuestos individuales. En una economía competitiva, en cambio, la factibilidad se logra a través de restricciones presupuestarias $p\cdot x_i \le p\cdot \omega_i + T_i$. Incluso si una asignación ($x^*,q^*$) es eficiente, puede no ser alcanzable desde dotaciones $\omega_i$ dadas, porque algunos consumidores no podrían financiar $x_i^*$ al precio $p$. Las transferencias de suma fija $T_i$ (o, equivalentemente, una redistribución de dotaciones) se introducen para garantizar viabilidad presupuestaria sin distorsionar decisiones marginales: al ser suma fija, no altera las CPO's, por lo que no rompe la eficiencia, pero sí hace que $x_i^*$ sea elegible para cada agente. 


En suma: primero, el planner identifica asignaciones eficientes y sus precios sombra; después, el problema con transferencias muestra que, con p igual a esos precios sombra y con una redistribución apropiada (transferencias), esa asignación eficiente puede ser un equilibrio competitivo.



\clearpage
\anexo{Modelo de la telaraña (Cobweb) – Información incompleta y ajuste lento}\label{anx:cobb_webb}
En el análisis seguido en la exposición hemos supuesto información perfecta en el sentido de que consumidores y empresas reaccionan simultánea e inmediatamente a los shocks.

Sin embargo, en algunos mercados los agentes no pueden conocer de inmediato los nuevos precios (o los conocen pero no actúan hasta saber si se trata de variaciones transitorias o permanentes del precio), por lo que el ajuste puede no ser instantáneo y los productores pueden necesitar más tiempo para ajustar su capacidad productiva a los nuevos precios (es el caso, por ejemplo, de la agricultura o de la construcción).\footnote{%
    Representa un caso de sluggishness para la reacción de los productores ante shocks. En el período actual la producción esta fija si hay un shock de reacciones en el período siguiente. Por tanto, la oferta no depende del precio actual sino del precio del período anterior. Esto conduce a oscilaciones del equilibrio.
} Así, en estos casos, la oferta no dependerá del precio en curso, sino del precio del período anterior.

Este modelo recibe el nombre de telaraña porque conduce a oscilaciones del equilibrio que se asemejan a la forma de una telaraña. Estas oscilaciones del equilibrio serán convergentes, divergentes o constantes en función de las elasticidades relativas de la oferta y la demanda.


\textsc{Fluctuación convergente}

En este caso, los precios tienden al equilibrio en el largo plazo. Si la demanda es más elástica que la oferta y partimos de un equilibrio como $P_t$. Para este nivel de precios, la cantidad ofertada aumentará hasta $Q_{t+1}$, pero para el nivel de oferta de $Q_{t+1}$, los consumidores sólo estarán dispuestos a pagar $P_{t+1}$. Para Pt+1 los productores solo ofertarán $Q_{t+2}$, etc. La fluctuaciones de precios se van reduciendo progresivamente y los precios y las cantidades tienden al equilibrio en el largo plazo. 


\imagenfit{Din_Convergente.png}{Dinámica convergente: Modelos lineales}{Apuntes Víctor Gutiérrez. Tema 3.A.16}

\textsc{Fluctuación divergente}

En este caso, los precios se desvían del equilibrio en el largo plazo. El razonamiento es inverso al anterior y esto sucede cuando la demanda es menos elástica que la oferta.

\imagenfit{Fluctuacion_Divergente.png}{Dinámica divergente: Modelos lineales}{Apuntes Víctor Gutiérrez. Tema 3.A.16}


\textsc{Fluctuación continua}

En caso de que la elasticidad de la demanda sea igual a la elasticidad de la oferta, los precios oscilarán alrededor del equilibrio sin alcanzarlo y sin divergir de él. Es decir, la magnitud de las oscilaciones de los precios será constante, produciendo un rectángulo.

\imagenfit{Estacionariedad.png}{Dinámica estable: Órbita sobre el \textit{fixed point} -- Modelos lineales}{Apuntes Víctor Gutiérrez. Tema 3.A.16}

\textsc{Fluctuación divergente hacia una fluctuación continua}

Hay otro escenario posible, que se da por los cambios en la pendiente de las curvas: si cerca del punto donde las curvas de oferta y demanda se cruzan la demanda es más pronunciada que la oferta, pero cuando nos movemos lejos de este punto la curva de oferta se hace más empinada que la curva de demanda, en un primer momento, los precios y las cantidades actúan como una fluctuación divergente. Sin embargo, como la oferta se hace más empinada que la demanda, un ciclo límite puede ser alcanzado, convirtiendo esta fluctuación divergente en una fluctuación continua.

\imagenfit{Divergencia_Estacionariedad.png}{Dinámica estable: Órbita sobre el \textit{fixed point} -- Cuasi-lineales}{Apuntes Víctor Gutiérrez. Tema 3.A.16}

\end{document}